\documentclass[12pt, a4paper]{ctexart}
\usepackage{geometry}
\usepackage{amsmath}
\usepackage{amssymb}
\usepackage{graphicx}
\usepackage{enumitem}
\usepackage{hyperref}

\geometry{left=2.5cm, right=2.5cm, top=2.5cm, bottom=2.5cm}

\title{冲鸭考研数学资料}
\author{公众号：冲鸭考研}
\date{2026年2月4日}

\begin{document}

\maketitle
\tableofcontents

\newpage

\section{第一章 函数、极限、连续}

\subsection{A类}

\subsubsection*{一、选择题}

\begin{enumerate}
    \item 在下列区间内，函数 $ f(x)=\frac{e^{3x}-1}{x(x-1)} $ 有界的是（）\\
    (A) $ \left(0,\frac{1}{2}\right) $. \quad (B) $ \left(\frac{1}{2},1\right) $. \quad (C) $ \left(1,+\infty\right) $. \quad (D) 以上都不正确.

    \item 在下列区间内，函数 $ f(x)=\frac{\left[1-\cos^{2}\left(x-4\right)\right]\sin\left(x-3\right)}{\left(x-4\right)^{2}\left(x-3\right)^{2}\left(x-2\right)} $ 有界的是（）\\
    (A) $ (1,2) $. \quad (B) $ (2,3) $. \quad (C) $ (3,4) $. \quad (D) $ (4,5) $.

    \item 下列命题中, 正确的是（）\\
    (A) 若 $ f(x) $ 与 $ g(x) $ 都不是连续函数，则 $ f(x) + g(x) $ 也不是连续函数.\\
    (B) 若 $ f(x) $ 与 $ g(x) $ 都不是单调函数，则 $ f(x) + g(x) $ 也不是单调函数.\\
    (C) 若 $ f(x) $ 与 $ g(x) $ 都是无界函数，则 $ f(x) + g(x) $ 也是无界函数.\\
    (D) 若 $ f(x) $ 与 $ g(x) $ 都是奇函数，则 $ f(x) + g(x) $ 也是奇函数.

    \item 设 $ \{a_{n}\}, \{b_{n}\} $ 均为正项数列，其中 $ \{a_{n}\} $ 无界， $ \{b_{n}\} $ 有界，则下列数列中，一定有界的是（）\\
    (A) $ \left\{\frac{b_{n}}{a_{n}}\right\}. $ \quad (B) $ \left\{\frac{b_{n}}{1+a_{n}b_{n}}\right\}. $ \quad (C) $ \left\{\sqrt[n]{a_{n}b_{n}}\right\} $ \quad (D) $ \left\{a_{n}b_{n}-a_{n-1}b_{n-1}\right\} $

    \item 设数列 $ \{a_n\} $ 满足 $ a_n = \sqrt[n]{n}, n = 1, 2, \cdots $，则下列命题中，正确的是（）\\
    (A) 数列 $ \{a_{n}\} $ 能取到最小值, 但取不到最大值.\\
    (B) 数列 $ \{a_{n}\} $ 能取到最大值, 但取不到最小值.\\
    (C) 数列 $ \{a_{n}\} $ 既能取到最大值, 又能取到最小值.\\
    (D) 数列 $ \{a_{n}\} $ 既不能取到最大值, 又不能取到最小值.

    \item 设 $ a_{n}=\left(1+\sin\frac{1}{n}\right)^{n}, n=1,2,\cdots $ ，则下列命题中，正确的是（）\\
    (A) 数列 $ \{a_{n}\} $ 单调增加且收敛于 e . \quad (B) 数列 $ \{a_{n}\} $ 单调减少且收敛于 e .\\
    (C) 数列 $ \{a_{n}\} $ 单调增加且收敛于 $ e^{\sin 1} $ . \quad (D) 数列 $ \{a_{n}\} $ 单调减少且收敛于 $ e^{\sin 1} $

    \item 设数列 $ \{x_{n}\} $ 收敛, 则（）\\
    (A) 当 $ \lim_{n\to\infty}\tan x_n=0 $ 时， $ \lim_{n\to\infty}x_n=0 $\\
    (B) 当 $ \lim_{n\to\infty}\left(x_{n}^{3}+\sqrt{\left|x_{n}\right|}\right)=0 $ 时， $ \lim_{n\to\infty}x_{n}=0 $\\
    (C) 当 $ \lim_{n\to\infty}\left(x_{n}^{2}+x_{n}^{3}\right)=0 $ 时， $ \lim_{n\to\infty}x_{n}=0 $\\
    (D) 当 $ \lim_{n\to\infty}\left[x_{n}+\ln(x_{n}+1)\right]=0 $ 时， $ \lim_{n\to\infty}x_{n}=0 $

    \item 设数列 $ \{a_{n}\} $ 满足 $ \lim_{n\to\infty}n(a_{n+1}-a_{n}-1)=0 $ ，则（）\\
    (A) $ \{a_{n}\} $ 有界.\\
    (B) $ \lim_{n\to\infty}a_n $ 存在.\\
    (C) $ \{a_{n}\} $ 自某项起单调减少.

    \item 设数列 $ \{a_{n}\}, \{b_{n}\} $ 满足 $ \lim_{n\to\infty}a_{n}=\infty,\lim_{n\to\infty}b_{n}=\infty $ ，则下列结论中，正确的是（）\\
    (A) $ \lim_{n\to\infty}e^{a_n}=\infty $.\\
    (B) $ \lim_{n\to\infty}\left(a_{n}+b_{n}\right)=\infty $.\\
    (C) $ \lim_{n\to\infty}\frac{a_n}{b_n}=1 $\\
    (D) $ \lim_{n\to\infty}a_n b_n=\infty $.

    \item 已知 $ \lim_{n\to\infty}a_n=a $ ，则下列关于数列 $ \{a_n\} $ 的说法中，正确的是（）\\
    (A) 若 a > 0，则当 n 充分大时， $ a_{n} > 0 $\\
    (B) 若 $ a \geq 0 $，则 $ \{a_{n}\} $ 中只存在有限个负项.\\
    (C) 若 a=0 ，则当 n 充分大时， $ a_{n}\geq-\frac{1}{n} $\\
    (D) 若 a=0 ，则当 n 充分大时， $ a_{n}\leq\frac{1}{n} $

    \item 当 $ x \to 1 $ 时，函数 $ \frac{2}{\mathrm{e}^{\frac{x}{x-1}} - 1} - \frac{\tan(x-1)}{|x-1|} $ 的极限（）\\
    (A) 等于 -1. \quad (B) 等于 1. \quad (C) 为 $ \infty $. \quad (D) 不存在但不为 $ \infty $.

    \item 设 $ \alpha_{1}=\sqrt{1+\tan x}-\sqrt{1+\sin x},\alpha_{2}=\int_{0}^{x^{4}}\frac{1}{\sqrt{1-t^{2}}}dt,\alpha_{3}=\int_{0}^{x^{2}}\sin t^{2}dt $．当 $ x\to0 $ 时，以上三个无穷小量按照从低阶到高阶的排序是（）\\
    (A) $ \alpha_{1},\alpha_{2},\alpha_{3} $ \quad (B) $ \alpha_{1},\alpha_{3},\alpha_{2} $ \quad (C) $ \alpha_{2},\alpha_{1},\alpha_{3} $ \quad (D) $ \alpha_{3},\alpha_{1},\alpha_{2} $

    \item 设 $ \alpha_{1}=\int_{0}^{x^{2}}\mathrm{d}u\int_{0}^{u}\arctan t\mathrm{d}t,\alpha_{2}=\int_{0}^{x}\mathrm{d}u\int_{0}^{u^{2}}\arctan t\mathrm{d}t,\alpha_{3}=\int_{0}^{x}\mathrm{d}u\int_{0}^{u}\arctan t^{2}\mathrm{d}t $．当 $ x\to0 $ 时，以上三个无穷小量按照从高阶到低阶的排序是（）\\
    (A) $ \alpha_{1},\alpha_{2},\alpha_{3} $ \quad (B) $ \alpha_{1},\alpha_{3},\alpha_{2} $ \quad (C) $ \alpha_{3},\alpha_{1},\alpha_{2} $ \quad (D) $ \alpha_{3},\alpha_{2},\alpha_{1} $

    \item 设 $ \alpha_1 = \sqrt{x + \sqrt{x}}, \alpha_2 = \sqrt[3]{x} \tan\left(x + \sqrt{x}\right), \alpha_3 = 1 - \cos\sqrt{x} $. 当 $ x \to 0^+ $ 时，以上三个无穷小量按照从低阶到高阶的排序是（ ）\\
    (A) $ \alpha_1, \alpha_2, \alpha_3 $. \quad (B) $ \alpha_1, \alpha_3, \alpha_2 $. \quad (C) $ \alpha_2, \alpha_1, \alpha_3 $. \quad (D) $ \alpha_3, \alpha_1, \alpha_2 $.

    \item 设 $ 1 - e^{x^2} = x \tan \alpha(x) $，其中 $ \left| \alpha(x) \right| < \frac{\pi}{2} $，则当 $ x \to 0 $ 时， $ \alpha(x) $ 是（）\\
    (A) 比 x 高阶的无穷小量.\\
    (B) 比 x 低阶的无穷小量.\\
    (C) 与 x 同阶但不等价的无穷小量.\\
    (D) 与 x 等价的无穷小量.

    \item 当 $ x \to 0 $ 时，若 $ \alpha_{1}(x), \alpha_{2}(x), \beta_{1}(x), \beta_{2}(x) $ 都是非零无穷小量，且 $ \alpha_{1}(x) \sim \alpha_{2}(x) $， $ \beta_{1}(x) \sim \beta_{2}(x) $，则下列命题中，错误的是（）\\
    (A) 若 $ \alpha_{1}(x) \sim \beta_{1}(x) $，则 $ \alpha_{2}(x) - \beta_{2}(x) = o(\alpha_{2}(x)) $。\\
    (B) 若 $ \alpha_{1}(x)-\beta_{1}(x)=o(\alpha_{1}(x)) $，则 $ \alpha_{2}(x)\sim\beta_{2}(x) $\\
    (C) 若 $ \alpha_{1}(x) \sim \beta_{1}(x) $，则 $ \alpha_{1}(x) - \beta_{1}(x) \sim \alpha_{2}(x) - \beta_{2}(x) $.\\
    (D) 若 $ \alpha_{1}(x)=o\left(\beta_{1}(x)\right) $，则 $ \alpha_{1}(x)-\beta_{1}(x)\sim\alpha_{2}(x)-\beta_{2}(x) $

    \item 当 $ x \to +\infty $ 时，无穷大量 $ I_{1} = x^{x^{x+1}} $, $ I_{2} = x^{(x+1)^{x}} $, $ I_{3} = (x+1)^{x^{x}} $ 按照从高阶到低阶的排序是（）\\
    (A) $ I_{1}, I_{2}, I_{3} $ \quad (B) $ I_{1}, I_{3}, I_{2} $ \quad (C) $ I_{2}, I_{1}, I_{3} $ \quad (D) $ I_{3}, I_{2}, I_{1} $

    \item 已知 $ \lim_{x\to0}\frac{(1-ax)^{\frac{1}{2}}+(1-bx)^{\frac{1}{4}}-2}{x^{2}}=-2 $ ，则 $ \left|a\right|+\left|b\right|= $ （）\\
    (A) 2 . \quad (B) 4 . \quad (C) 6 . \quad (D) 8 .

    \item 若 $ \lim_{n\to\infty}\frac{n^{a}}{\left(n+1\right)^{b}-n^{b}}=900 $ ，则（）\\
    (A) $ a = -\frac{899}{900}, b = \frac{1}{900} $. \quad (B) $ a=\frac{899}{900}, b=-\frac{1}{900} $\\
    (C) $ a=\frac{899}{900}, b=\frac{1}{900} $ \quad (D) $ a = -\frac{899}{900}, b = -\frac{1}{900} $

    \item 函数 $ f(x)=\frac{|x|^{2x}-1}{x(x+2)\ln|x|} $ 的可去间断点的个数为（）\\
    (A) 0 . \quad (B) 1 . \quad (C) 2. \quad (D) 3 .

    \item 下列关于函数 $ f(x)=\frac{|x|^{|x|}-x}{|x|^{|x|}\ln|x|} $ 的说法中, 正确的是（）\\
    (A) $ f(x) $ 只有可去间断点.\\
    (B) $ f(x) $ 只有跳跃间断点.\\
    (C) $ f(x) $ 有可去间断点和无穷间断点, 没有跳跃间断点.\\
    (D) $ f(x) $ 有可去间断点、跳跃间断点和无穷间断点.

    \item 函数 $ f(x)=\lim_{t\to0}\left(1+\frac{\sin t}{x}\right)^{\frac{\sin x}{t}} $ 在 $ (-\infty,+\infty) $ 内（）\\
    (A) 连续. \quad (B) 有可去间断点. \quad (C) 有跳跃间断点. \quad (D) 有无穷间断点.

    \item 设函数 $ f(x)=\frac{1}{\mathrm{e}^{\frac{(x-2)(x-3)}{x-1}}-1} $，则（）\\
    (A) x=1, x=2, x=3 都是 $ f(x) $ 的第一类间断点.\\
    (B) x=1, x=2, x=3 都是 $ f(x) $ 的第二类间断点.\\
    (C) x=1 是 $ f(x) $ 的第一类间断点，x=2, x=3 是 $ f(x) $ 的第二类间断点.\\
    (D) x=1 是 $ f(x) $ 的第二类间断点，x=2, x=3 是 $ f(x) $ 的第一类间断点.

    \item 设函数 $ f(x) = \lim_{n \to \infty} \frac{\left(e^{nx} + 1\right) \sin x}{1 + x^3 e^{nx}} $，则下列结论中，正确的是（）\\
    (A) $ f(x) $ 不存在间断点.\\
    (B) x = 0 是 $ f(x) $ 的可去间断点.\\
    (C) x=0 是 $ f(x) $ 的跳跃间断点.\\
    (D) x=0 是 $ f(x) $ 的无穷间断点.

    \item 函数 $ f(x)=\frac{(x-1)\sin x\arctan\frac{1}{x-1}}{\left(x^{4}-1\right)(x-2)|x|} $ 的跳跃间断点的个数是（）\\
    (A) 1 . \quad (B) 2 . \quad (C) 3 . \quad (D) 4 .

    \item 设函数 $ f(x) = \begin{cases} -1, & x < 0, \\ 1, & x \geq 0, \end{cases} \quad g(x) = \begin{cases} 2 + ax^2, & x \leq -1, \\ x^2, & -1 < x < 0, \\ x^2 - b, & x \geq 0. \end{cases} $ 若 $ f(x) + g(x) $ 在 R 上连续，则（）\\
    (A) a=1,b=1. \quad (B) a=1,b=2. \quad (C) a=-1,b=1. \quad (D) a=-1,b=2.

    \item 若函数 $ f(x)=\begin{cases}\frac{a e^{\frac{1}{x}}+b e^{-\frac{1}{x}}}{e^{\frac{2}{x}}-e^{-\frac{1}{x}}},&x\neq0,\\0,&x=0\end{cases} $ 在 R 上连续，则（）\\
    (A) $ a=0, b=0 $ \quad (B) $ a \neq 0, b = 0 $ \quad (C) a 可以取任意值，b = 0 \quad (D) a = 0, b 可以取任意值.

    \item 设函数 $ f(x) = \begin{cases} a\cos x, & x \geq 0, \\ x + 1, & x < 0, \end{cases} \quad g(x) = \begin{cases} x^2 - b^2, & x \geq 0, \\ x + \pi, & x < 0. \end{cases} $ 已知 $ b \neq 0 $，若 $ f(x) $ 连续，且 $ f(g(x)) $ 在 x = 0 处连续，则 $ a^2 + b^2 = $ （ ）\\
    (A) 1 . \quad (B) 2 . \quad (C) 3 . \quad (D) 5 .
\end{enumerate}

\subsubsection*{二、填空题}

\begin{enumerate}[start=29]
    \item $ \lim_{x \to 0} \frac{1 - x \cot x}{\sin x \tan x} = $ \underline{\hspace{2cm}}.
    
    \item $ \lim_{x \to +\infty} \frac{\sqrt[3]{x^3 + x^2} - \sqrt{x}}{\ln(e^x + x^3)} = $ \underline{\hspace{2cm}}.

    \item $ \lim_{x\to0}\frac{\left(\cos x-\sqrt{\cos x}\right)\sin\left(\sin x\right)}{\left[x-\ln\left(1+\tan x\right)\right]\left(\mathrm{e}^{x}-1\right)}= $ \underline{\hspace{2cm}}.

    \item $ \lim_{x\to\infty}\left(\frac{x^{2}+x+4}{x^{2}+2x+3}\right)^{x^{2}\sin\frac{2}{x}}= $ \underline{\hspace{2cm}}.
    
    \item $ \lim_{x\to0}(\cos x)^{\frac{1}{x-\ln(1+x)}}= $ \underline{\hspace{2cm}}.

    \item 设函数 $ f(x)=\int_{0}^{x}\sin x e^{t^{2}}dt $，则 $ \lim_{x\to0}\left(2-\cos x\right)^{\frac{1}{f(x)}}= $ \underline{\hspace{2cm}}.

    \item $ \lim_{x\to1^{-}}\left(\sqrt{\frac{1}{1-x}+100}-\sqrt{\frac{1}{1-x}-100}\right)= $ \underline{\hspace{2cm}}.
    
    \item $ \lim_{x\to0}\left[\frac{1}{1-\cos2x}-\frac{1}{2\ln\left(1+x^{2}\right)}\right]= $ \underline{\hspace{2cm}}.

    \item $ \lim_{x \to 0^+} \left( \frac{1}{e^{\sqrt{x}} - 1} \right)^{\frac{1}{\ln x}} = $ \underline{\hspace{2cm}}.

    \item $ \lim_{x \to +\infty} \left(1 + \frac{1}{x}\right)^{x^3} e^{-x^2 + \frac{x}{2}} = $ \underline{\hspace{2cm}}.
    
    \item $ \lim_{x \to +\infty} \left[ \sin \ln(1 + x) - \sin \ln x \right] = $ \underline{\hspace{2cm}}.

    \item 已知当 $ x \to 1 $ 时， $ x^{100} - 1 $ 是 $ x^{kx} - 1 $ 的等价无穷小，则 $ k = $ \underline{\hspace{2cm}}.

    \item 已知当 $ x \to 0^{+} $ 时， $ e^{x} + 2\cos\sqrt{x} - 3 $ 是 $ 1 - \cos ax $ 的等价无穷小，则 $ a^{2} = $ \underline{\hspace{2cm}}.
    
    \item 当 $ x \to \infty $ 时， $ \left[\frac{e}{\left(1+\frac{1}{x}\right)^x}\right]^x - \sqrt{e} $ 与 $ x^k $ 是同阶无穷小量，则 $ k = $ \underline{\hspace{2cm}}.

    \item 已知函数 $ f(x) $ 连续，且 $ f(0)=1 $，则 $ \lim_{x\to0}\frac{1-\cos\left[xf(x)\right]}{\left(\sqrt{1+x^2}-1\right)f(x)}= $ \underline{\hspace{2cm}}.

    \item 已知 $ \lim_{x\to0}\frac{\ln\left[1+\frac{f(x)}{\sin2x}\right]}{e^x-1}=1 $ ，则 $ \lim_{x\to0}\frac{f(x)}{1-\cos x}= $ \underline{\hspace{2cm}}.
    
    \item 已知 $ \lim_{x\to0}\left(a\left[x\right]+\frac{4b}{\pi}\arctan\frac{1}{x}+\frac{4+e^{-\frac{1}{x}}}{1+e^{-\frac{2}{x}}}\right)=1 $ ，其中 $ [x] $ 表示不超过 $ x $ 的最大整数，则 $ ab= $ \underline{\hspace{2cm}}.

    \item 设 k 为整数，函数 $ f(x)=\frac{kx}{k+1+e^{kx}} $ 在 $ (-\infty,+\infty) $ 上连续，且 $ \lim_{x\to-\infty}f(x) $ 存在，则 k= \underline{\hspace{2cm}}.

    \item 设函数 $ f(x) = \frac{x}{x-1} $，当 $ x \neq 1 $ 时， $ f_n(x) = f\{f\{\cdots\{f\{f(x)\}\}\} $，即 $ n $ 个 $ f(x) $ 复合所得函数，则 $ \lim_{x \to 0} \frac{f_{2n+1}(x)}{f_{2n}(x)} = $ \underline{\hspace{2cm}}.
    
    \item $ \lim_{n\to\infty}\left[\frac{1}{2!}+\frac{2}{3!}+\cdots+\frac{n}{\left(n+1\right)!}\right]^{2n\cdot n!}= $ \underline{\hspace{2cm}}.

    \item $ \lim_{n\to\infty}\sqrt[n]{\left|\cos0\right|1^{2}+\left|\cos1\right|2^{2}+\cdots+\left|\cos\left(n-1\right)\right|n^{2}}= $ \underline{\hspace{2cm}}.

    \item $ \lim_{n\to\infty}\left(\frac{1}{3n^{2}+2\sin^{2}1+1}+\frac{3}{3n^{2}+2\sin^{2}2+2}+\cdots+\frac{2n-1}{3n^{2}+2\sin^{2}n+n}\right)= $ \underline{\hspace{2cm}}.
    
    \item $ \lim_{n\to\infty}\frac{1!+2!+\cdots+n!}{n\cdot n!}= $ \underline{\hspace{2cm}}.
\end{enumerate}

\subsubsection*{三、解答题}

\begin{enumerate}[start=52]
    \item 设函数 $ f(x) $ 以 2 为周期. 记 $ a_n = f(n) $. 若 $ \lim_{n \to \infty} a_n $ 存在, 证明: 数列 $ \{a_n\} $ 为常数列.

    \item 求极限 $ \lim_{x\to0}\frac{x-\sin x}{x-x\tan^{2}x-\tan x} $

    \item 求极限 $ \lim_{x\to0}\frac{x-\sin\left(xe^{-x}\right)}{\sin x\ln\left(x+1\right)} $

    \item 求极限 $ \lim_{x\to0}\frac{1-\cos x\cos2x\cos3x}{x^{2}} $

    \item 已知函数 $ f(x) $ 具有一阶连续导数且 $ f(0) \neq 0 $，求极限 $ \lim_{x \to 0} \left[ \frac{1}{\int_{0}^{x^{2}} f(t) dt} - \frac{1}{x^{2} f\left(x^{2}\right)} \right] $.

    \item 已知当 $ x \to 0 $ 时， $ ax^{2} + bx^{4} + c \int_{0}^{x} \frac{\sin t^{2}}{t} dt $ 与 $ x^{6} $ 是等价无穷小，求 $ a + b + c $。

    \item 已知极限 $ \lim_{x\to0}\frac{\cos(a\sin x)-b\cos x}{x^{4}} $ 存在，求 a,b 的值并计算该极限.

    \item 已知当 $ x \to 0^{+} $ 时， $ \arcsin^{\alpha}x, (x - \sin x)^{\frac{1}{\alpha}} $ 均为比 x 高阶的无穷小量，且函数 $ f(x) = \begin{cases} \frac{\arcsin^{\alpha}x}{x^{\beta}(x - \sin x)^{\frac{1}{\alpha}}}, & x > 0, \\ \gamma, & x \leq 0. \end{cases} $ 若 $ \alpha $ 为整数，且 $ f(x) $ 连续，求 $ \alpha\beta\gamma $.

    \item 设函数 $ f(x) $ 为 $ (-\infty, +\infty) $ 上连续的奇函数，且满足 $ f(a - x) = f(a + x) $，其中 $ a > 0 $。证明：\\
    (I) $ f(x) $ 为周期函数;\\
    (II) $ f(x) $ 在 $ (-\infty, +\infty) $ 上有界.

    \item 设 $ ab \neq 0 $，证明方程 $ x = a \sin x + b \cos x + \sqrt{a^2 + b^2} $ 至少有一个正根，并且不超过 $ 2\sqrt{a^2 + b^2} $。\\
    (I) 证明方程 $ e^{-ax} + e^{-(n-1)x} + \cdots + e^{-x} = 1 (n \geq 2) $ 在区间 $ (0,1) $ 内有且仅有一个实根.\\
    (II) 记(I) 中的实根为 $ x_{n} $，证明 $ \lim_{n\to\infty}x_{n} $ 存在，并求此极限.
\end{enumerate}

\subsection{B 类}

\subsubsection*{一、选择题}

\begin{enumerate}
    \item 设 $ f(x) $ 为 $ [0,\pi] $ 上的连续函数, 则下列命题中, 正确的是（）\\
    (A) 若 $ f(\operatorname{arccot} x) $ 在 $ [0, \pi] $ 上单调增加，则 $ f(x) $ 在 $ [0, \pi] $ 上单调增加.\\
    (B) 若 $ f(\operatorname{arccot} x) $ 在 $ [0, \pi] $ 上单调增加，则 $ f(x) $ 在 $ [0, \pi] $ 上单调减少.\\
    (C) 若 $ f(x) $ 在 $ [0, \pi] $ 上单调增加，则 $ f(\operatorname{arccot} x) $ 在 $ [0, \pi] $ 上单调增加.\\
    (D) 若 $ f(x) $ 在 $ [0, \pi] $ 上单调增加，则 $ f(\operatorname{arccot} x) $ 在 $ [0, \pi] $ 上单调减少.

    \item 已知数列 $ \{a_{n}\} $ 和连续函数 $ f(x) $ 。若 $ b_{n}=f(a_{n}) $ ，且 $ \lim_{n\to\infty}b_{n}=1 $ ，则下列条件中，可以推出“数列 $ \{a_{n}\} $ 收敛”的条件个数是（）\\
    ① $ f(x) $ 是单调函数. ② $ f(x) $ 是偶函数. ③ $ f(x) $ 是有界函数. ④ $ f(x) $ 是周期函数.\\
    (A) 0 . \quad (B) 1 . \quad (C) 2. \quad (D) 3 .

    \item 函数 $ f(x)=x-[x]+|\sin x| $, $ x\in(-\infty,+\infty) $ 是 （ ）\\
    (A) 周期函数. (B) 单调函数. (C) 奇函数. (D) 有界函数.

    \item 设函数 $ f(x) $ 满足 $ f(x+4) - f(x) = f(2) $，则下列结论中，正确的个数为（）\\
    ① 若 $ f(x) $ 为奇函数，则 $ f(x) $ 为周期函数.\\
    ② 若 $ f(x) $ 为偶函数, 则 $ f(x) $ 为周期函数.\\
    ③ 若 $ f(x) $ 为周期为 4 的函数，则 $ f(x) $ 为奇函数.\\
    ④ 若 $ f(x) $ 为周期为 4 的函数，则 $ f(x) $ 为偶函数.\\
    (A) 1 . \quad (B) 2 . \quad (C) 3 . \quad (D) 4 .

    \item 下列数列中，无界的是（）\\
    (A) $ x_{1} $ 为任意实数， $ x_{n+1}=\frac{x_{n}}{1+x^{2}} $\\
    (B) $ x_{1} $ 为任意实数， $ x_{n+1}=\left(3x_{n}^{2}-2x_{n}+1\right)\mathrm{e}^{-x_{n}^{2}} $\\
    (C) $ x_{1} $ 为任意大于 2 的实数， $ x_{n+1}=\frac{1}{2}x_{n}+1 $。\\
    (D) $ x_{1} $ 为任意正实数， $ x_{n+1}=x_{n}(x_{n}+1) $

    \item 设函数 $ f(x)=\frac{(x-1)\sin x}{x+1}\sqrt{1+\frac{1}{x^{2}}} $，则下列关于数列 $ \left\{f(n)\right\},\left\{f\left(\frac{1}{n}\right)\right\}(n=1,2,3,\cdots) $ 的说法中，正确的是（）\\
    (A) $ \left\{f(n)\right\},\left\{f\left(\frac{1}{n}\right)\right\} $ 均有界.\\
    (B) $ \left\{f(n)\right\} $ 有界, $ \left\{f\left(\frac{1}{n}\right)\right\} $ 无界.\\
    (C) $ \left\{f(n)\right\} $ 无界, $ \left\{f\left(\frac{1}{n}\right)\right\} $ 有界. \\
    (D) $ \left\{f(n)\right\},\left\{f\left(\frac{1}{n}\right)\right\} $ 均无界.

    \item 已知数列 $ \{x_{n}\}, \{y_{n}\} $ 满足 $ x_{1}=y_{1}=\frac{1}{2}, x_{n+1}=1-\cos x_{n}, y_{n+1}=y_{n}^{2}(n=1,2,\cdots) $，则当 $ n \to \infty $ 时，（）\\
    (A) $ \lim_{n\to\infty}\frac{x_n-y_n}{y_n}=0 $\\
    (B) $ \lim_{n\to\infty}\frac{x_n+y_n}{y_n}=1 $.\\
    (C) $ \lim_{n\to\infty}\frac{x_n-y_n}{x_n}=0 $.\\
    (D) $ \lim_{n\to\infty}\frac{x_n+y_n}{x_n}=1 $.

    \item 已知数列 $ \{x_{n}\} $，其中 $ -\infty < x_{n} < +\infty $，函数 $ f(x) = e^{-|x|} $，则（）\\
    (A) 当 $ \lim_{n\to\infty}\arctan f(x_n) $ 存在时， $ \lim_{n\to\infty}x_n $ 存在.\\
    (B) 当 $ \lim_{n\to\infty}f\left(\arctan x_{n}\right) $ 存在时， $ \lim_{n\to\infty}x_{n} $ 存在.\\
    (C) 当 $ \lim_{n\to\infty}\arctan f(x_n) $ 存在时， $ \lim_{n\to\infty}f(x_n) $ 存在，但 $ \lim_{n\to\infty}x_n $ 不一定存在.\\
    (D) 当 $ \lim_{n\to\infty}f\left(\arctan x_{n}\right) $ 存在时， $ \lim_{n\to\infty}\arctan x_{n} $ 存在，但 $ \lim_{n\to\infty}x_{n} $ 不一定存在.

    \item 设数列 $ \{x_{n}\} $ 与 $ \{y_{n}\} $ 满足 $ \lim_{n\to\infty}x_{n}y_{n}=1 $ ，则下列结论中，正确的是（）\\
    (A) 若 $ \{x_{n}\} $ 发散，则 $ \{y_{n}\} $ 必收敛.\\
    (B) 若 $ \left\{x_{n}\right\} $ 无界，则 $ \left\{y_{n}\right\} $ 必有界.\\
    (C) 若 $ \{x_{n}\} $ 有界，则 $ \{y_{n}\} $ 必收敛.\\
    (D) 若 $ \left\{\frac{1}{x_{n}}\right\} $ 为无穷小，则 $ \{y_{n}\} $ 必为无穷小.

    \item 已知数列 $ \{a_{n}\},\{b_{n}\} $，满足 $ 0<a_{n}<b_{n}<2a_{n} $，则下列命题中，正确的是（）\\
    ① 若 $ \lim_{n\to\infty}a_{n}=0 $ ，则 $ \lim_{n\to\infty}b_{n}=0 $\\
    ② 若 $ \lim_{n\to\infty}b_{n}=0 $ ，则 $ \lim_{n\to\infty}a_{n}=0 $\\
    ③ 若 $ \lim_{n\to\infty}a_{n} $ 存在, 则 $ \lim_{n\to\infty}b_{n} $ 存在.\\
    ④ 若 $ \lim_{n\to\infty}\left(a_{n}-b_{n}\right) $ 存在，则 $ \lim_{n\to\infty}a_{n}\leq\lim_{n\to\infty}b_{n} $\\
    (A) ①②. \quad (B) ②③. \quad (C) ①②③. \quad (D) ①③④.

    \item 已知 $ a_{n}=\frac{\sin n}{n}(n=1,2,\cdots) $，则数列 $ \{a_{n}\} $（）\\
    (A) 有最大值，有最小值.\\
    (B) 有最大值，没有最小值.\\
    (C) 没有最大值，有最小值.\\
    (D) 没有最大值，没有最小值.

    \item 若 $ \lim_{x\to0}\frac{\ln\left(1+3x^{2}\right)+xf(x)}{x^{4}}=\frac{1}{2} $，则 $ \lim_{x\to0}\frac{3x+f(x)}{x^{3}}= $（）\\
    (A) $ \frac{1}{2} $. (B) 1. (C) 3. (D) 5.

    \item 已知 $ a_n = \lim_{x \to 1} \frac{(1-x)(1-\sqrt{2x-1})\cdots[1-\sqrt{nx-(n-1)}]}{n(x-1)^n} $，则（ ）\\
    (A) $ \lim_{n \to \infty} a_n, \lim_{n \to \infty} na_n $ 均存在.\\
    (B) $ \lim_{n \to \infty} a_n $ 存在， $ \lim_{n \to \infty} na_n $ 不存在.\\
    (C) $ \lim_{n \to \infty} na_n, \lim_{n \to \infty} |na_n| $ 均存在.\\
    (D) $ \lim_{n \to \infty} na_n $ 存在， $ \lim_{n \to \infty} |na_n| $ 不存在.

    \item 已知函数：① $ f(x)=2x $；② $ f(x)=x^{2} $；③ $ f(x)=\sqrt{|x|} $；④ $ f(x)=\arctan x $，则满足对于任意的首项 $ a_{1} $，数列 $ a_{n}=f(a_{n-1})(n=2,3,\cdots) $ 均收敛的函数个数为（）\\
    (A) 1. \quad (B) 2. \quad (C) 3. \quad (D) 4.

    \item 设函数 $ f(x) $ 满足 $ \int_{0}^{x} f(t) \, \mathrm{d}t + 2f(x) + f'(x) = \mathrm{e}^{-x} \sin 2x $，且 $ f(0) = f'(0) = 0 $，则当 $ x \to 0 $ 时，与 $ f(x) $ 为同阶但不等价的无穷小量为（ ）\\
    (A) $ \sin 2x $ . \quad (B) $ 1 - \cos 2x $ . \quad (C) $ e^{x^2} - 1 $ . \quad (D) $ \sqrt{1 + x} - 1 $ .

    \item 设函数 $ f(x) $ 在 R 上处处有定义，且 $ f(0) = 0 $，则下列命题中，错误的是（）\\
    (A) 当 $ x \to 0 $ 时, 若 $ f(x) \sim \sin^{2} x $, 则 $ f'(0) $ 存在.\\
    (B) 若 $ 0 \leq f(x) \leq \sin^{2} x $ 恒成立，则 $ f'(0) $ 存在.\\
    (C) 若在 $ [0,+\infty) $ 上 $ g(x) \leq f(x) \leq h(x) $，在 $ (-\infty,0) $ 上 $ h(x) \leq f(x) \leq g(x) $，且当 $ x \to 0 $ 时，函数 $ g(x) $ 和 $ h(x) $ 都是 x 的同阶无穷小，则 $ f(x) $ 也是 x 的同阶无穷小。\\
    (D) 当 $ x \to 0 $ 时, 若 $ f'(0) $ 存在且不为 0, 则 $ f(x) $ 是 x 的同阶无穷小.

    \item 设函数 $ f(x) $ 连续， $ F(x) $ 为 $ f(x) $ 的原函数，则下列命题中，错误的是（）\\
    (A) 当 $ x \to 0 $ 时, 若 $ f(x) = o(1) $, 则 $ F(x) = o(x) $.\\
    (B) 当 $ x \to 0 $ 时, 若 $ F(x) = o(x) $, 则 $ f(x) = o(1) $.\\
    (C) 当 $ x \to 0 $ 时, 若 $ f(x) = o(1) $ , 则 $ \int_{0}^{x} f(t) \, dt = o(x) $.\\
    (D) 当 $ x \to 0 $ 时, 若 $ \int_{0}^{x} f(t) dt = o(x) $，则 $ f(x) = o(1) $.

    \item 设函数 $ f(x)=\lim_{n\to\infty}\arctan\frac{1}{1+\sin^n x} $，则下列关于 $ f(x) $ 的命题中，正确的是（）\\
    ① $ f(x) $ 是周期为 $ 2\pi $ 的奇函数. ② $ f(x) $ 有无穷多个可去间断点.\\
    ③ $ f(x) $ 有无穷多个跳跃间断点. ④ $ f(x) $ 的最小值小于 $ \frac{\pi}{6} $.\\
    (A) ①②. \quad (B) ③④. \quad (C) ①④. \quad (D) ②④.

    \item 函数 $ f(x)=\frac{\left(a-\arctan\frac{1}{x}\right)\left(\mathrm{e}^{x}-b\right)}{x^{2}(x-b)} $ ( $ b\neq0 $) 的可去间断点的个数最多为（）\\
    (A) 0 . \quad (B) 1 . \quad (C) 2 . \quad (D) 3 .
\end{enumerate}

\subsubsection*{二、填空题}

\begin{enumerate}[start=20]
    \item 设函数 $ f(x) = \lim_{n \to \infty} \frac{\ln\left(\left|\cos x\right| + \arctan^{2n-1}x\right)}{\sqrt{n}} $，记 $ f(x) $ 的自然定义域为 $ (a, b] $，则 $ b - a = $ \underline{\hspace{2cm}}.
    
    \item 设函数 $ f(x)=2x^{3}+x^{2}+1, x>0, x=f^{-1}(y) $ 为 $ f(x) $ 的反函数，若正整数 k 满足 $ \lim_{t\to+\infty}\frac{f^{-1}(2t)}{k\sqrt{t}} $ 存在且非零，则 k= \underline{\hspace{2cm}}.

    \item 已知常数 a > 0, b > 0，若 $ \lim_{x \to 0} \left( \frac{a^x + b^x}{2b^x} \right)^{\frac{1}{x}} = \sqrt{b}, \lim_{x \to 0} \left( \frac{1 + bx}{1 - bx} \right)^{\frac{b}{2ax}} = \frac{a}{b} $，则 a = \underline{\hspace{2cm}}.

    \item 设 $ a > 0, \lim_{x \to \infty} \frac{x \sqrt{a e^{2|x|} + e^{|x|}} - a(x + \ln|x|) e^{|x|}}{\sqrt{x^2 + \ln|x|} e^{|x|}} $ 存在，则 $ a = $ \underline{\hspace{2cm}}.

    \item 设当 0 < x ≤ 1 时，函数 $ f(x) = x^{\sin x} $，对于其他 $ x, f(x) = 2f(x + 1) - \ln k $。若 $ f(x) $ 在 x = 0 处连续，则 k = \underline{\hspace{2cm}}.

    \item 设函数 $ f(x)=\ln\frac{1-x}{1+x} $，若在 x=0 的邻域内存在函数 $ \xi=\xi(x) $，使得 $ f(x)=x^{2}f'(\xi) $，则 $ \lim_{x\to0}\frac{\xi^{2}-1}{x}= $ \underline{\hspace{2cm}}.

    \item 设数列 $ \{a_{n}\} $ 满足 $ a_{1}=1, a_{2}=2 $ ，且 $ a_{n}=\frac{1}{3}a_{n-1}+\frac{2}{3}a_{n-2}(n=3,4,5,\cdots) $ ，则 $ \lim_{n\to\infty}a_n= $ \underline{\hspace{2cm}}.

    \item 设 $ a_{n}=\int_{1}^{+\infty}x^{-3}\ln^{n}x\mathrm{d}x\left(n=0,1,2,\cdots\right),b_{n}=\frac{a_{n}}{(2n)!!} $，则 $ b_{n}= $ \underline{\hspace{2cm}}.

    \item $ \lim_{n\to\infty}\left(\arctan\frac{2}{1+0+0^{2}}+\arctan\frac{2}{1+2+1^{2}}+\cdots+\arctan\frac{2}{1+2n+n^{2}}\right)= $ \underline{\hspace{2cm}}.

    \item 设 $ a_{n}=\int_{\frac{n}{n+1}}^{\frac{2n}{2n+1}}\frac{x^{n-1}}{1+x^{n}}dx $，则 $ \lim_{n\to\infty}na_{n}= $ \underline{\hspace{2cm}}.

    \item 设区域 $ D=\left\{(x,y)\middle|\arcsin y\leq x\leq\frac{\pi}{4},0\leq y\leq\frac{\sqrt{2}}{2}\right\} $，对正整数 $ n,f_{n}(x,y)=\frac{ny^{n-1}}{x+e^{x}} $，则 $ \lim_{n\to\infty}\iint_{D}f_{n}(x,y)\mathrm{d}x\mathrm{d}y= $ \underline{\hspace{2cm}}.
\end{enumerate}

\subsubsection*{三、解答题}

\begin{enumerate}[start=31]
    \item 计算 $ \lim_{x\to+\infty}\sqrt{x}\sin\left(\sqrt{x+2\sqrt{x}}-\sqrt{x+\sqrt{x}}-\frac{1}{2}\right) $

    \item (I) 求 $ \lim_{x\to+\infty}\frac{\arctan2x-\arctan x}{\frac{\pi}{2}-\arctan x} $;\\
    (II) 若 $ \lim_{x\to+\infty}x\left[1-f(x)\right] $ 不存在，而 $ I=\lim_{x\to+\infty}\frac{\arctan2x+\left[b-1-bf(x)\right]\arctan x}{\frac{\pi}{2}-\arctan x} $ 存在，试确定 b 的值，并求 I.

    \item 已知 a, b 为正整数，若极限 $ I_{1} = \lim_{x \to +\infty} \frac{bx^{b}}{x^{2a} - (x + 1)^{a}(x - 1)^{a}} $ 和 $ I_{2} = \lim_{x \to 0^{+}} \frac{bx^{b}}{x^{2a} - (1 + x)^{a}(1 - x)^{a} + 1} $ 等于同一个非零常数，求 a, b 的值.

    \item 设函数 $ f(x) $ 满足 $ f(0) = f'(0) = 0 $，且 $ \lim_{x \to 0} \frac{a \sin x - \frac{x}{1 + bx^2}}{f(x)} = 1 $， $ \lim_{x \to 0} \frac{ax - \frac{\sin x}{1 + bx^2}}{f(x)} = 2 $ 。求 a, b 的值.

    \item 设函数 $ f(x) $ 在 $ [0, +\infty) $ 上具有一阶连续导数， $ f(0) = 0 $，且 $ \lim_{x \to +\infty} \left[ f(x) + f'(x) \right] = a > 0 $。\\
    (I) 记 $ g(x) = \mathrm{e}^{x} f(x) $，求 $ g'(x) $;\\
    (II) 证明：存在 $ x_{0}>0 $ ，使得当 $ x>x_{0} $ 时， $ f(x)>\frac{a}{2} $

    \item 设数列 $ \left\{x_{n}\right\} $ 满足 $ x_{n+1}=\sqrt{\frac{\pi}{2}}x_{n}\sin x_{n} $，且 $ 0<x_{1}<\frac{\pi}{2} $．求 $ \lim_{n\to\infty}\frac{\sec x_{n}-\tan x_{n}}{\frac{\pi}{2}-x_{n}} $．

    \item 计算 $ \lim_{n\to\infty}\frac{n\left[\ln^{3}\left(n+2\right)-\ln^{3}n\right]}{\ln^{2}\left(n+1\right)} $

    \item 已知 $ f(x) $ 为 $ (0,+\infty) $ 上连续的有界函数，且满足 $ \lim_{x\to0^{+}}\frac{f(x)}{x-\tan x}=6 $ 。证明：方程 $ f(x)+x^{3}=0 $ 在 $ (0,+\infty) $ 上必有一实根。

    \item 设函数 $ f(x) $ 在区间 $ [0,900] $ 上连续，且 $ f(0) + 900 = f(900) $，证明：存在 $ x_0 \in [0,900] $，使得 $ f(x_0) + 1 = f(x_0 + 1) $。

    \item 设 $ x_{n+1}=\frac{1}{3}\left(2x_{n}+\frac{27}{x_{n}^{2}}\right) $，其中 $ x_{1}>0 $ 。证明数列 $ \{x_{n}\} $ 收敛，并求 $ \lim_{n\to\infty}x_{n} $ 。

    \item 设数列 $ \{x_{n}\} $ 满足 $ 0 < x_{1} < 1, x_{n+1} = \ln(1 + x_{n}) (n = 1, 2, \cdots) $.\\
    (I) 证明 $ \lim_{n\to\infty}x_n $ 存在，并求该极限；\\
    (II) 计算 $ \lim_{n\to\infty}\left(\frac{x_{n+1}}{x_n}\right)^{\frac{1}{x_n}} $

    \item 已知数列 $ \{x_n\} $ 满足 $ x_n = \ln(2n-1)(n=1,2,\cdots) $，前 $ n $ 项和为 $ X_n $， $ \{y_n\} $ 满足 $ y_n = \ln 2n(n=1,2,\cdots) $，前 $ n $ 项和为 $ Y_n $。若 $ s_n = X_n - Y_n $，证明：\\
    (I) 数列 $ \left\{e^{s_{n}}\right\} $ 收敛；\\
    (II) $ -\frac{1}{2}\ln4n < s_n < -\frac{1}{2}\ln2n $.

    \item 设函数 $ f(x) $ 在 $ (-\infty, +\infty) $ 内有定义，且在 x = 0 处连续。对任意的 x，均有 $ f(x) = f(1 - \cos x) $。\\
    (I) 设数列 $ \left\{x_{n}\right\} $ 满足 $ x_{0}=c_{0} $ 且 $ x_{n+1}=1-\cos x_{n}(n=0,1,2,\cdots) $. 证明： $ \lim_{n\to\infty}x_n $ 存在，并求 $ \lim_{n\to\infty}x_n $.\\
    (II) 证明： $ f(x) $ 在 $ (-\infty, +\infty) $ 内恒为常数.

    \item （I）比较 $ \int_{0}^{1}\left(\frac{\sin x}{2x}\right)^{n}dx $ 与 $ \int_{0}^{1}\left(1-\frac{x}{2}\right)^{n}dx\left(n=1,2,\cdots\right) $ 的大小，说明理由；\\
    (II) 记 $ a_{n}=\int_{0}^{1}\left(\frac{\sin x}{2x}\right)^{n}\mathrm{d}x\left(n=1,2,\cdots\right) $，求极限 $ \lim_{n\to\infty}a_{n} $

    \item （I）证明：当 x>0 时， $ \frac{x}{1+x}<\ln(1+x)<x $\\
    （II）设 $ x_{n}=\left(1+\frac{n^{2}-n+1}{n^{3}}\right)\left(1+\frac{n^{2}-n+3}{n^{3}}\right)\cdots\left(1+\frac{n^{2}+n-1}{n^{3}}\right) $，求 $ \lim_{n\to\infty}x_n $。
\end{enumerate}

\subsection{c 类}

\subsubsection*{一、选择题}

\begin{enumerate}
    \item 设 $ f(x), g(x) $ 均为 $ [0, +\infty) $ 上的连续函数，若对任意正数 x，有 $ f(2x) = f(x), g(x^2) = g(x) $，且 $ f(0) > g(0) $，则（）\\
    (A) 对任意 $ x \geq 0 $ ，均有 $ f(x) > g(x) $\\
    (B) 对任意 $ x \geq 0 $ ，均有 $ f(x) = g(x) $\\
    (C) 对任意 $ x \geq 0 $ ，均有 $ f(x) < g(x) $\\
    (D) 存在 $ x_{1} \geq 0, x_{2} \geq 0 $ ，使得 $ f(x_{1}) > g(x_{1}), f(x_{2}) < g(x_{2}) $

    \item 设数列 $ \{x_{n}\} $ 和 $ \{y_{n}\} $ 满足 $ \begin{cases} x_{1}=1, \\ y_{1}=1 \end{cases} $ 和 $ \begin{cases} x_{n+1}=y_{n}, \\ y_{n+1}=\frac{x_{n}}{2}, \end{cases} $ 则下列说法中，错误的是（）\\
    (A) $ \lim_{n\to\infty}x_n=0 $\\
    (B) $ \lim_{n\to\infty}y_n=0 $\\
    (C) $ \lim_{n\to\infty}(x_n+y_n)=0 $\\
    (D) 以上说法均不正确.
\end{enumerate}

\subsubsection*{二、解答题}

\begin{enumerate}[start=3]
    \item 求极限 $ \lim_{n\to\infty}\frac{1}{n}\prod_{k=1}^{n}\frac{k+1-\sqrt{k^{2}+k}}{\sqrt{k}\left(\sqrt{k+2}-\sqrt{k+1}\right)} $

    \item 设周期函数 $ f(x) = (-1)^{[x]} \left| x - \left[ x + \frac{1}{2} \right] \right| $，其中 $ [x], \left[ x + \frac{1}{2} \right] $ 分别表示不超过 $ x, x + \frac{1}{2} $ 的最大整数，记 $ a_n = \int_0^1 \frac{f(nx)}{x} \, dx $ 。证明：\\
    (Ⅰ) 数列 $ \{a_{2n-1}\} $ 单调减少， $ \{a_{2n}\} $ 单调增加；\\
    (II) $ \lim_{n\to\infty}a_{n} $ 存在.
\end{enumerate}

\section{第二章 一元函数微分学}

\subsection{A类}

\subsubsection*{一、选择题}

\begin{enumerate}
    \item 设函数 $ f(x) $ 可导， $ y = f(x) \cos x $ 当自变量 x 在点 x = 0 处取得增量 $ \Delta x = 0.2 $ 时，相应的函数增量 $ \Delta y $ 的线性主部为 0.4，则 $ f'(0) = (\quad) $\\
    (A) 2. \quad (B) 0.2. \quad (C) 1. \quad (D) 0.1.

    \item 设函数 $ \varphi(x)=\begin{cases}x^{2}\left(2+\sin\frac{1}{x}\right),&x\neq0,\\0,&x=0,\end{cases} $ 且函数 $ f(x) $ 在 x=0 处可导，则函数 $ f(\varphi(x)) $ 在 x=0 处（）\\
    (A) 不连续.\\
    (B) 连续但不可导.\\
    (C) 可导且导数为 0 .\\
    (D) 可导且导数不为 0 .

    \item 设函数 $ f(x) = \begin{cases} \arctan \frac{x+1}{x-1} + a, & x > 1, \\ c, & x = 1, \\ \arctan \frac{x+1}{x-1} + b, & x < 1 \end{cases} $ 可导，则 $ f'(1) = (\quad) $\\
    (A) $ -\frac{1}{2} $. (B) $ \frac{1}{2} $. (C)1. (D)与a,b的值有关.

    \item 若函数 $ f(x), g(x) $ 满足 $ f(x) = f(-x), g(x) = -g(-x) $，在 $ (0, +\infty) $ 内， $ f'(x) > 0 $， $ f''(x) > 0 $， $ g''(x) > 0 $，则在 $ (-\infty, 0) $ 内恒成立的是（）\\
    (A) $ f'(x) > g'(x) $.\\
    (B) $ f'(x) < g'(x) $.\\
    (C) $ f''(x) > g''(x) $.\\
    (D) $ f''(x) < g''(x) $.

    \item 设函数 $ f(x) = \sin x \cdot \ln x \cdot \ln 2x \cdots \ln nx $，其中 $ n $ 为正整数，则 $ f'(1) = $ ( )\\
    (A) $ \sin 1 \cdot \ln 2 \ln 3 \cdots \ln n $ . (B) $ \ln 2 \ln 3 \cdots \ln n $ . (C) 1 . (D) 不存在.

    \item 设函数 $ f(x) = (1 - \cos x)(2 - \cos x) $... ( $ n - \cos x $)，则 $ f''(0) = $ （ ）\\
    (A) $ (n - 1)! $ . \quad (B) $ n! $ . \quad (C) $ (n + 1)! $ . \quad (D) 0 .

    \item 若函数 $ f(x)=\begin{cases}\lim\limits_{t\to-\infty}\dfrac{\sin\left[tx\left(\sqrt{t^{2}+2e}+t\right)\right]}{x\arctan t},&x\neq0,\\\dfrac{2e}{\pi},&x=0,\end{cases} $ 则下列说法中，正确的是（）\\
    (A) $ f(x) $ 在 x=0 处的极限存在,但不连续.\\
    (B) $ f(x) $ 在 x=0 处连续,但不可导.\\
    (C) $ f(x) $ 在 x = 0 处可导，且 $ f'(0) = 0 $.\\
    (D) $ f(x) $ 在 x = 0 处可导，且 $ f'(0) \neq 0 $.

    \item 设函数 $ f(x)=\frac{1}{2x^{3}-5x^{2}+4x} $，则下列说法中，正确的是（）\\
    (A) $ f(x) $ 恰有一个单调增加区间, 两个单调减少区间.\\
    (B) $ f(x) $ 恰有一个单调增加区间, 三个单调减少区间.\\
    (C) $ f(x) $ 恰有一个单调减少区间, 两个单调增加区间.\\
    (D) $ f(x) $ 恰有一个单调减少区间, 三个单调增加区间.

    \item 设 $ f(x) $ 为 $ (-\infty, +\infty) $ 上的单调增加函数，则下列说法中，正确的是（）\\
    (A) $ \left[f(x)\right]^{2} $ 单调增加.\\
    (B) $ \left[-f(-x)\right]^{3} $ 单调增加.\\
    (C) 若 $ f(x) $ 可导，则对任意 $ x \in (-\infty, +\infty) $，都有 $ f'(x) > 0 $。\\
    (D) 若 $ f(x) $ 可导，则 $ f(x) $ 存在反函数 $ x = f^{-1}(y) $，且 $ f^{-1}(y) $ 可导.

    \item 设函数 $ f(x)=\begin{cases}-x^{3}\ln x,&x>0,\\0,&x=0,\\\arctan\left(x+\frac{1}{x}\right)+\frac{\pi}{2},&x<0,\end{cases} $ 则（）\\
    (A) $ f(x) $ 有两个极大值点, 无极小值点. (B) $ f(x) $ 有一个极大值点, 一个极小值点.\\
    (C) $ f(x) $ 有两个极大值点, 一个极小值点. (D) $ f(x) $ 有一个极大值点, 两个极小值点.

    \item 下列四个命题中, 正确的是（）\\
    (A) 设 $ x_0 \in (a, b) $，函数 $ f(x) $ 在 $ (a, x_0) $ 内满足 $ f'(x) < 0 $，在 $ (x_0, b) $ 内满足 $ f'(x) > 0 $，则 $ f(x) $ 在 $ x = x_0 $ 处取得它在 $ (a, b) $ 上的最小值。\\
    (B) 设函数 $ f(x) $ 在 $ x = x_{0} $ 处取得极小值，则存在 $ \delta > 0 $，使 $ f(x) $ 在 $ (x_{0} - \delta, x_{0}) $ 内单调减少，在 $ (x_{0}, x_{0} + \delta) $ 内单调增加.\\
    (C) 设 $ f(x) $ 为区间 $ (-a, a) $ 上的偶函数，则 x=0 为 $ f(x) $ 的一个极值点.\\
    (D) 设 $ f(x) $ 在区间 $ (-a, a) $ 内二阶可导且为奇函数，则 $ f''(0) = 0 $

    \item 设函数 $ f(x) $ 在 $ (-\infty, +\infty) $ 内有定义， $ x_{0} \neq 0 $ 是函数 $ f(x) $ 的极小值点，则（）\\
    (A) $ x_{0} $ 必是 $ f(x) $ 的驻点.\\
    (B) $ -x_{0} $ 必是 $ f(-x) $ 的极大值点.\\
    (C) $ -x_{0} $ 必是 $ -f(-x) $ 的极大值点.\\
    (D) $ -x_{0} $ 必是 $ -f(x) $ 的极大值点.

    \item 设函数 $ f(x) $ 在 x = a 的某个邻域内具有二阶导数，且 $ f''(a) \neq 0 $，且存在 $ \delta > 0 $，使得当 $ x \in (a - \delta, a + \delta) $ 时， $ \lim_{t \to a} \frac{f(t) - f(x)}{(t - x)^2} $ 均存在，且当 $ x \neq a $ 时， $ \lim_{t \to a} \frac{f(t) - f(x)}{(t - x)^2} \geq 0 $，则（）\\
    (A) x=a 是 $ f(x) $ 的极小值点.\\
    (B) x = a 是 $ f(x) $ 的极大值点.\\
    (C) x = a 不是 $ f(x) $ 的极值点.\\
    (D) 无法确定 x = a 是否为极值点.

    \item 已知函数 $ f(x) $ 在 x=0 的某个邻域内连续，且 $ f(0)=0,\lim_{x\to0}\frac{f(x)}{1-\mathrm{e}^{x^2}}=1 $ ，则 $ f(x) $ 在 x=0 处（）\\
    (A) 不可导. (B) 可导, 且 $ f'(0) \neq 0 $ . (C) 取得极大值. (D) 取得极小值.

    \item 设函数 $ f(x) $ 和 $ g(x) $ 在 $ [0,1] $ 上的最大值分别为 a, b，则下列关于 $ f(x)g(x) $ 在 $ [0,1] $ 上的最大值 c 的说法中，正确的是（）\\
    (A) c < ab . (B) c > ab . (C) c = ab . (D) 以上都有可能.

    \item 已知函数 $ f(x) $ 在定义域内可导, 它的图形如右图所示, 则导函数 $ f'(x) $ 的图形可能为 ( )\\
    (A) [图] (B) [图] (C) [图] (D) [图]

    \item 设函数 $ f(x) $ 在 $ (-\infty, +\infty) $ 上连续, 其导函数 $ f'(x) $ 的图形如右图所示, 则下列说法中, 正确的是（）\\
    (A) 函数 $ f(x) $ 有 1 个极值点，曲线 $ y = f(x) $ 有 1 个拐点.\\
    (B) 函数 $ f(x) $ 有 2 个极值点，曲线 $ y = f(x) $ 有 2 个拐点.\\
    (C) 函数 $ f(x) $ 有 3 个极值点，曲线 $ y = f(x) $ 有 2 个拐点.\\
    (D) 函数 $ f(x) $ 有 3 个极值点，曲线 $ y = f(x) $ 有 1 个拐点.

    \item 设 $ f(x) $ 是 $ (0,+\infty) $ 上单调增加的正值函数, 且存在二阶导数, 则对于满足 $ 0<x_{1}<x_{2} $ 的任意 $ x_{1}, x_{2} $ , 下列说法中, 正确的是（）\\
    (A) 若 $ f''(x) > 0 $，则 $ \sqrt{f(x_{1})f(x_{2})} < f\left(\frac{x_{1} + x_{2}}{2}\right) $.\\
    (B) 若 $ f''(x) > 0 $，则 $ f\left(\sqrt{x_1 x_2}\right) < \frac{f\left(x_1\right) + f\left(x_2\right)}{2} $.\\
    (C) 若 $ f''(x)<0 $ ，则 $ \sqrt{f(x_{1})f(x_{2})}>f\left(\frac{x_{1}+x_{2}}{2}\right) $\\
    (D) 若 $ f''(x)<0 $ ，则 $ f\left(\sqrt{x_{1}x_{2}}\right)>\frac{f\left(x_{1}\right)+f\left(x_{2}\right)}{2} $

    \item 曲线 $ y=(x-1)x^{\frac{4}{3}} $ 的拐点个数为（）\\
    (A) 0 . \quad (B) 1 . \quad (C) 2 . \quad (D) 3 .

    \item 下列曲线中有斜渐近线的是（）\\
    (A) $ y = \ln x + \cos \frac{1}{x} $ \quad (B) $ y = x + \cos \frac{1}{x} $\\
    (C) $ y = \ln \frac{1}{x} + \cos x $ \quad (D) $ y = \frac{1}{x} + \cos x $

    \item 曲线 $ y = xe^{\frac{1}{x}} $ 的渐近线的条数为（）\\
    (A) 0 . \quad (B) 1 . \quad (C) 2 . \quad (D) 3 .

    \item 曲线 $ y=\sqrt{2x^{2}+4x+1} $ 的渐近线的条数为（）\\
    (A) 1 . \quad (B) 2 . \quad (C) 3 . \quad (D) 4 .

    \item 曲线 $ y = x \ln \frac{x}{x-1} + \ln \left[ x(x-1) \right] $ 的渐近线的条数为（）\\
    (A) 0 . \quad (B) 1 . \quad (C) 2. \quad (D) 3 .

    \item 设函数 $ y(x)=\lim_{t\to0}\left[1-\frac{\ln(1-t)}{x^{2}}\right]^{\frac{x}{\sin t}} $ ，则下列关于曲线 $ y=y(x) $ 的渐近线的说法中，正确的是（）\\
    ① 该曲线无渐近线. ② 该曲线有铅直渐近线.\\
    ③ 该曲线有水平渐近线. ④ 该曲线有斜渐近线.\\
    (A) ②. \quad (B) ③. \quad (C) ②③. \quad (D) ②④.

    \item 设函数 $ f(x)=\frac{x^{2}+x+a}{x+b} $, $ g(x)=\frac{x+\sqrt{x^{2}+4(1-b)x+a}}{2} $，则当 $ x\to+\infty $ 时，下列关于曲线 $ y=f(x) $ 与 $ y=g(x) $ 的渐近线的结论中，正确的是（）\\
    (A) 没有公共斜渐近线.\\
    (B) 是否具有公共斜渐近线与 a, b 的取值有关.\\
    (C) 有公共斜渐近线, 且公共斜渐近线的方程与 a 的取值有关.\\
    (D) 有公共斜渐近线, 且公共斜渐近线的方程与 b 的取值有关.

    \item 设函数 $ f(x)=\ln^{2}x, g(x)=\frac{1}{x}, h(x)=\cot^{2}x $ ，则当 $ x \to 0^{+} $ 时，有（）\\
    (A) $ g(x) < h(x) < f(x) $ \quad (B) $ h(x) < g(x) < f(x) $.\\
    (C) $ f(x) < g(x) < h(x) $ \quad (D) $ g(x) < f(x) < h(x) $

    \item 设函数 $ f(x) = \begin{cases} x^{4} \sin \frac{1}{x}, & x \neq 0, \\ 0, & x = 0, \end{cases} $ 则使得 $ f^{(n)}(x) $ 连续的最高阶数 n 为（）\\
    (A) 0 . \quad (B) 1 . \quad (C) 2. \quad (D) 3 .

    \item 设函数 $ f(x) = \begin{cases} x^{\alpha} \ln \frac{1}{x}, & x > 0, \\ 0, & x = 0, \\ e^{-\frac{1}{x^2}}, & x < 0. \end{cases} $ 若 $ f'(x) $ 在 $ x = 0 $ 处连续，则（）\\
    (A) $ 0 < \alpha < 1 $ \quad (B) $ \alpha > 1 $ \quad (C) $ 0 < \alpha < 2 $ \quad (D) $ \alpha > 2 $

    \item 设函数 $ y = f(x) $ 由 $ \begin{cases} x = t |t| \\ y = |t| \sin t^2 \end{cases} $ 确定，则（）\\
    (A) $ f(x) $ 连续， $ f'(0) $ 不存在.\\
    (B) $ f'(0) $ 存在， $ f'(x) $ 在 x=0 处不连续.\\
    (C) $ f'(x) $ 连续， $ f''(0) $ 不存在.\\
    (D) $ f''(0) $ 存在， $ f''(x) $ 在 x=0 处不连续.

    \item 设函数 $ f(x) = \begin{cases} e^{-\frac{1}{x}}, & x > 0, \\ |\tan x|, & -\frac{\pi}{2} < x \leq 0, \end{cases} $ 则 $ x = 0 $ 是 $ f(x) $ 的（）\\
    (A) 可导点，极值点. (B) 不可导点，极值点. (C) 可导点，非极值点. (D) 不可导点，非极值点.

    \item 设函数 $ f(x) = (x - 1) \left| \ln x \cdot \ln \left( x + \frac{1}{2} \right) \right| $，则下列关于 $ f(x) $ 的说法中，正确的是（）\\
    (A) x=1 不是极值点, 也不是驻点.\\
    (B) x=1 是极值点,但不是驻点.\\
    (C) $ x = \frac{1}{2} $ 既是极值点, 又是驻点.\\
    (D) $ x = \frac{1}{2} $ 是极值点, 但不是驻点.

    \item 设函数 $ f(x) $ 二阶可导, 满足 $ f'(a) = 0 $ 且 $ \lim_{x \to a} \frac{f''(x)}{x - a} = 1 $，则（）\\
    (A) x=a 是 $ f(x) $ 的极小值点.\\
    (B) x=a 是 $ f(x) $ 的极大值点.\\
    (C) $ (a, f(a)) $ 是曲线 $ y = f(x) $ 的拐点.\\
    (D) x = a 不是 $ f(x) $ 的极值点， $ (a, f(a)) $ 也不是曲线 $ y = f(x) $ 的拐点.

    \item 设函数 $ f(x) $ 在 x=0 附近可导，且 $ \lim_{x\to0}\frac{f(x)}{x} $ 存在，则下列命题中，错误的是（）\\
    (A) $ f(0)=0 $\\
    (B) $ f'(0) $ 存在.\\
    (C) $ \lim_{x\to0}f'(x)=0 $\\
    (D) 在 x=0 附近，存在常数 C>0 ，使得 $ \left|f(x)\right|<C|x| $

    \item 已知方程 $ x^{4}-2x^{2}-k=0 $ 有四个不同的实根，则 k 的取值范围是（）\\
    (A) $ \left(-\infty,0\right) $. \quad (B) $ \left(-1,0\right) $. \quad (C) $ \left(-1,1\right) $. \quad (D) $ \left(1,+\infty\right) $.

    \item 设对于任意 $ \alpha \in \left(0, \frac{\pi}{2}\right) $，方程 $ x^{\cos^2 \alpha} = k + x \cos^2 \alpha (x > 0) $ 有两个不同的实根，则 k 的取值范围是（）\\
    (A) $ \left[0,\sin^{2}\alpha\right) $. \quad (B) $ \left(0,\sin^{2}\alpha\right) $. \quad (C) $ \left[0,\cos^{2}\alpha\right) $. \quad (D) $ \left(0,\cos^{2}\alpha\right) $.

    \item 设函数 $ f(x)=\cot 2x $ 在 $ x=\frac{\pi}{4} $ 处的 3 次泰勒多项式为 $ a+b\left(x-\frac{\pi}{4}\right)+c\left(x-\frac{\pi}{4}\right)^2+ d\left(x-\frac{\pi}{4}\right)^3 $，则 $ ac+bd= $ （ ）\\
    (A) $ -\frac{8}{3} $ \quad (B) $ \frac{8}{3} $ \quad (C) $ -\frac{16}{3} $ \quad (D) $ \frac{16}{3} $

    \item 设函数 $ f(x)=\frac{\tan x}{1-x^{2}} $ 在 x=0 处的 3 次泰勒多项式为 $ ax+bx^{2}+cx^{3} $，则（）\\
    (A) $ a=1, b=0, c=-\frac{2}{3} $ \quad (B) $ a=1, b=0, c=\frac{2}{3} $\\
    (C) $ a=1, b=0, c=-\frac{4}{3} $ \quad (D) $ a=1, b=0, c=\frac{4}{3} $

    \item 设函数 $ f(x) $ 满足 $ f'(x) > f(x) $，则（）\\
    (A) $ ef(0) < f(1) < e^{2} f(-1) $\\
    (B) $ e^{2}f(-1) < f(1) < ef(0) $.\\
    (C) $ \mathrm{ef}(0) < \mathrm{e}^{2} f(-1) < f(1) $.\\
    (D) $ e^{2}f(-1) < ef(0) < f(1) $.
\end{enumerate}

\subsubsection*{二、填空题}

\begin{enumerate}[start=39]
    \item 设函数 $ y = f(\tan x)(\tan x)^n $，其中 $ f(x) $ 可微， $ f(1) = f'(1) = 1 $，则 $ y' \left( \frac{\pi}{4} \right) = $ \underline{\hspace{2cm}}.

    \item 设正值函数 $ f(x) $ 可微， $ g(x)=\arctan(2f(x)+1) $， $ f'(0)=1 $， $ g'(0)=\frac{1}{2} $，则 $ f(0)= $ \underline{\hspace{2cm}}.

    \item 设函数 $ y = f\left(\ln\sqrt{1 + x^2}\right) $，其中 $ f(x) $ 可微， $ f'(1) = 1 $， $ f''(1) = 0 $，则 $ y''\left(\sqrt{e^2 - 1}\right) = $ \underline{\hspace{2cm}}.

    \item 设函数 $ f(x) = \begin{cases} \ln x, & x > 1, \\ \sqrt{3} \arcsin \frac{x}{2} - \frac{\sqrt{3}\pi}{6}, & x \leq 1, \end{cases} y = f\left(f(x)\right) $，则 $ y'(e) = $ \underline{\hspace{2cm}}.

    \item 设函数 $ y = y(x) $ 由方程 $ e^{x^2 + y^2} + \sin(xy) = 0 $ 确定，则 $ y'(x) = $ \underline{\hspace{2cm}}.

    \item 设函数 $ y = y(x) $ 由方程 $ x + y - \frac{e^{x} \sin y}{3} = 0 $ 确定，则 $ y''(0) = $ \underline{\hspace{2cm}}.

    \item 设函数 $ f(x)=\int_{a}^{x}\ln\left(1+\csc^{2}t\right)\mathrm{d}t\left(0<a<\frac{\pi}{2}\right) $，若 $ y=f(x) $ 的反函数 $ x=f^{-1}(y) $ 在 y=0 处的导数 $ \left.\frac{\mathrm{d}x}{\mathrm{d}y}\right|_{y=0}=\frac{1}{\ln5} $，则 $ a= $ \underline{\hspace{2cm}}.

    \item 设函数 $ y = y(x) $ 由参数方程 $ \begin{cases} x = f(\arcsin 2t) + 1, \\ y = f(\sin t + \cos t - 1) \end{cases} $ 确定，其中 $ f(u) $ 可微，且 $ f'(0) \neq 0 $，则 $ \left. \frac{dy}{dx} \right|_{t=0} = $ \underline{\hspace{2cm}}.

    \item 设函数 $ y = y(x) $ 由参数方程 $ \begin{cases} x = t^2 e^t, \\ y = \arctan t^2 \end{cases} $ 确定，则 $ \left. \frac{d^2 y}{dx^2} \right|_{t=1} = $ \underline{\hspace{2cm}}.

    \item 设函数 $ f(x)=\ln\left(\sqrt{1+\sin^{2}x}+\sin x\right) $，则 $ f^{(4)}(0)= $ \underline{\hspace{2cm}}.

    \item $ \left(\frac{x^{4}+1}{x+1}\right)^{(5)}= $ \underline{\hspace{2cm}}.

    \item 设函数 $ f(x)=x^{2}\cos2x $ ，则 $ f^{(10)}(0)= $ \underline{\hspace{2cm}}.

    \item 设函数 $ f(x)=\mathrm{e}^{x}\cos x+\mathrm{e}^{-x}\sin x $，则 $ f^{(99)}(0)= $ \underline{\hspace{2cm}}.

    \item 曲线 $ \begin{cases} x = t^{2}, \\ y = t^{3} - 3t \end{cases} $ 在 $ t = \sqrt{3} $ 对应点处的切线方程为 \underline{\hspace{2cm}}.

    \item 曲线 $ \begin{cases} x = e^{-\frac{1}{t^2}}, \\ y = \arccos t \end{cases} $ 在 $ t = \frac{1}{2} $ 对应点处的法线方程为 \underline{\hspace{2cm}}.

    \item 设曲线 L 的参数方程为 $ \begin{cases} x = \ln \tan \frac{t}{2} + \cos t, \\ y = \sin t, \end{cases} $ 其中 $ \frac{\pi}{2} < t < \pi $，则曲线上一点 M 处的切线与 x 轴的交点 P 与点 M 之间的距离为 \underline{\hspace{2cm}}.

    \item 设曲线 L 的极坐标方程是 $ r\theta=2 $，则 L 在点 $ (r,\theta)=\left(\frac{2}{\pi},\pi\right) $ 处的切线的直角坐标方程是 \underline{\hspace{2cm}}.

    \item 曲线 $ y = x^{2} + 2x - 2 (x > 0) $ 上曲率半径为 $ 3\sqrt{6} $ 的点的纵坐标为 \underline{\hspace{2cm}}.

    \item 曲线 $ y + xy - e^{x} + e^{y} = 0 $ 在点 $ (0, y(0)) $ 处的曲率为 \underline{\hspace{2cm}}.

    \item 曲线 $ \begin{cases} x = t - \sin t, \\ y = 1 - \cos t \end{cases} $ 在 $ t = \frac{\pi}{2} $ 对应点处的曲率半径为 \underline{\hspace{2cm}}.

    \item 已知动点 $ P(x, y) $ 在曲线 $ y = x^{3} $ 上运动, 过点 $ P $ 作 $ x $ 轴的垂线, 交 $ x $ 轴于点 $ A $, 过点 $ P $ 作 $ y $ 轴的垂线, 交 $ y $ 轴于点 $ B $. 记矩形 $ OAPB $ 的面积为 $ S $. 若动点 $ P $ 的横坐标对时间的变化率为常数 $ v_{0} $, 则当点 $ P $ 运动到点 $ (1,1) $ 时, $ S $ 对时间的变化率为 \underline{\hspace{2cm}}.

    \item 如图所示, 在距离海岸边点 A 180m 的点 P 处有一灯塔, 灯塔上的探照灯光以每秒 r 弧度 (rad) 的速率匀速转动, 探照光线与 PA 的所成夹角为 $ \theta $. 若灯光经过岸边距离点 A 60m 处的点 B 时, 速率为 30 m/s, 则 $ r = $ \underline{\hspace{2cm}}.

    \item 记函数 $ f_{n}(x)=x^{n}\mathrm{e}^{-3x}(n>0,x>0) $ 的极值为 $ a_{n} $，则数列 $ \{a_{n}\} $ 的最小值为 \underline{\hspace{2cm}}.

    \item 当 $ x > e^{-1} $ 时，曲线 $ y = 2x \ln x - 5x + 3 $ 在 $ x = $ \underline{\hspace{2cm}} 处取得最小曲率半径。

    \item 曲线 $ y = \sin x e^{-x} $ 在区间 $ (0, \pi) $ 内的拐点坐标为 \underline{\hspace{2cm}}.

    \item 曲线 $ y=\frac{2x^{5}+x^{3}-1}{x^{4}+1}+e^{\arctan x^{2}} $ 的斜渐近线方程为 \underline{\hspace{2cm}}.
\end{enumerate}

\subsubsection*{三、解答题}

\begin{enumerate}[start=65]
    \item 设函数 $ y(x)=(1+x)^{x}(x>0) $，记 $ y=y(x) $ 的反函数为 $ x=x(y) $，求 $ \left.\frac{dx}{dy}\right|_{y=2},\left.\frac{d^{2}x}{dy^{2}}\right|_{y=2} $.

    \item 设函数 $ y = y(x) $ 由方程组 $ \begin{cases} x = e^{t-1} + t^{3} - 2, \\ y + y^{3} \tan \frac{\pi}{4} t = -2(-2 < t < 2) \end{cases} $ 确定，求 $ \left.\frac{dy}{dx}\right|_{x=0} $

    \item 设函数 $ f(x) $ 在 x = 0 处可导，在 x = 0 的某个邻域内满足等式 $ 2f(x) + f(-x) = 3 + \frac{\sqrt{1 + \tan x} - \sqrt{1 + \sin x}}{x^2} $，且 $ f(x) $ 是周期为 5 的周期函数，求 $ y = f(x) $ 在点 $ (5, f(5)) $ 处的切线方程.

    \item 已知函数 $ f(x)=\begin{cases}\sqrt{x-\sin x}, & 0\leq x<\pi, \\ x\ln(-x), & x<0,\end{cases} $ 求 $ f'(x) $，并求 $ f(x) $ 的极值.

    \item 设函数 $ y(x) $ 由方程 $ 27y^{3} + 2x^{3} - 3x^{2} - 12x + 20 = 0 $ 所确定，求 $ y = y(x) $ 的极值点，并判断它们是否为驻点.

    \item 设函数 $ f(x) $ 在 $ x_0 $ 处连续且可导，在 $ x_0 $ 的某个去心邻域 $ \hat{U}(x_0) $ 内二阶可导，且 $ f'(x_0) = 0 $， $ \lim_{x \to x_0} f''(x) > 0 $ 。证明： $ x = x_0 $ 为 $ f(x) $ 的极小值点。

    \item 在底面半径为 1，高为 2 的圆锥体中挖去一个圆柱体，当被挖圆柱体的体积 V 最大时，圆柱体的底面半径和高分别为多少，并求出此最大体积.

    \item 在椭圆 $ \frac{x^{2}}{a^{2}}+\frac{y^{2}}{b^{2}}=1(a>0,b>0,a\neq b) $ 位于第一象限的部分上求一点 P，使得该点处的法线与两坐标轴所围成的三角形面积最大.

    \item 求函数 $ f(x)=xe^{\frac{1}{x}} $ 的单调区间与凹、凸区间.

    \item 设函数 $ f(x) $ 在 $ [-1,1] $ 上连续，且当 $ x \neq 0 $ 时， $ f(x) = \frac{e^{ax} + e^{bx} - 2}{\cos x + \cos 2x - 2} $，其中 a, b 为非零常数. 求 $ f(0) $，并讨论 $ f(x) $ 在 x = 0 处的可导性.

    \item 设函数 $ f(x)=\begin{cases}x^{a}\sin\frac{1}{x},&x>0,\\b,&x=0,\\\frac{1-\cos x}{(-x)^{a-2}},&x<0\end{cases} $，有连续的导函数，求 a 的取值范围。

    \item 设当 x > 0 时，方程 $ kx^{2} + \frac{1}{x} = 1 $ 有两个不同的解，求 k 的取值范围.

    \item 设函数 $ f(x)=\begin{cases}x^{\ln x}, & x>0, \\ x^{2}e^{\frac{1}{x}}, & x<0,\end{cases} $ 根据 k 的不同取值，讨论方程 $ f(x)=k $ 的实根情况.

    \item 讨论方程 $ e^{\cos^{2}x}\sin x = k $ 在开区间 $ \left(0, \frac{\pi}{2}\right) $ 内的根的个数，其中 k 为参数.

    \item 讨论方程 $ \left(x-\frac{4}{3}\right)\mathrm{e}^{\frac{6}{x^{2}}}=k $ 的不同实根的个数, 其中 k 为参数.

    \item 证明：当 -1 < x < 1 时， $ e^{x^{2}} + \ln\left(1 - x^{2}\right)\left(\frac{1 + x}{1 - x}\right)^{x}\geq 1 + 2x^{2} $.

    \item 证明：当 x > 0 时， $ \frac{1}{3}x^{3} + \frac{1}{2}x^{2} - x + \frac{1}{6} \geq x^{2} \ln x $.
\end{enumerate}

\subsection{B 类}

\subsubsection*{一、选择题}

\begin{enumerate}
    \item 设函数 $ f(x) $ 为 $ (-\infty, +\infty) $ 上的单调增加的正值函数，且 $ f'(0) $ 存在，则下列关于极限 $ I_{1} = \lim_{x \to 0^+} \left[ \frac{f(x)}{f(0)} \right]^{\frac{1}{x}} $ 和 $ I_{2} = \lim_{x \to 0^+} \left( \frac{1}{x} \right)^{f(x) - f(0)} $ 的结论中，正确的是（）\\
    (A) $ I_{1} $ 与 $ I_{2} $ 是否存在均与 $ f(x) $ 的表达式有关.\\
    (B) $ I_{1} $ 与 $ I_{2} $ 均存在，且 $ I_{1} > I_{2} $\\
    (C) $ I_{1} $ 与 $ I_{2} $ 均存在，且 $ I_{1}<I_{2} $\\
    (D) 以上说法均不正确.

    \item 设函数 $ f(x) $ 满足 $ f(0) = 0 $，则 $ f(x) $ 在 x = 0 处可导的充分必要条件为（）\\
    (A) $ \lim_{h\to0}\frac{1}{h^{3}}f(\tanh-h) $ 存在.\\
    (B) $ \lim_{h\to0}\frac{1}{h^{2}}f\left(\ln(1+h)-h\right) $\\
    (C) $ \lim_{h\to0}\frac{1}{h}f\left(\arctan h-h\right) $ 存在.\\
    (D) $ \lim_{h\to0}\frac{1}{h}\left[f(h)-f(-h)\right] $ 存在.

    \item 设函数 $ f(x) $ 在 x=0 处连续, 则下列命题中, 错误的是（）\\
    (A) 若 $ \lim_{x\to0}\frac{f(x)}{x}=a $ ，则 $ f'(0)=a $\\
    (B) 若 $ f(0)=0, f'(0)=a $ ，则 $ \lim_{x\to0}\frac{f(x)}{x}=a $\\
    (C) 若 $ \lim_{x\to0}\frac{f(x)}{x^{2}}=\frac{a}{2} $，则 $ f''(0)=a $\\
    (D) 若 $ f(0)=f'(0)=0, f''(0)=a $ ，则 $ \lim_{x\to0}\frac{f(x)}{x^{2}}=\frac{a}{2} $

    \item 设函数 $ f(x) $ 具有二阶导数， $ \Delta x $ 为自变量 x 在 $ x_{0} $ 处的增量， $ \Delta y $ 与 dy 分别为 $ f(x) $ 在 $ x_{0} $ 处对应的增量与微分，则下列说法中，正确的是（）\\
    (A) 若在 $ x_{0} $ 的某邻域内有 $ \mathrm{d}y < \Delta y $，则 $ f''(x) $ 在该邻域内恒大于 0。\\
    (B) 若在 $ x_{0} $ 的某邻域内有 $ |dy| < |\Delta y| $，则 $ f''(x) $ 在该邻域内恒大于 0。\\
    (C) 若 $ f''(x) $ 在 $ x_{0} $ 的某邻域内恒大于 0，则在该邻域内有 $ dy < \Delta y $\\
    (D) 若 $ f''(x) $ 在 $ x_0 $ 的某邻域内恒大于 0，则在该邻域内有 $ |dy| < |\Delta y| $

    \item 设函数 $ f(x) $ 具有二阶导数， $ \Delta x $ 为自变量 $ x $ 在 $ x_0 $ 处的增量， $ \Delta y $ 与 $ \mathrm{d}y $ 分别为 $ f(x) $ 在 $ x_0 $ 处对应的增量与微分，则 $ \mathrm{d}y < \Delta y $ 是 $ f''(x) $ 在 $ x_0 $ 的某邻域内大于 0 的（）\\
    (A) 充分必要条件.\\
    (C) 必要不充分条件.

    \item 设函数 $ f_1(x) $, $ f_2(x) $, $ g(x) $ 在 $ x=0 $ 的某邻域内均可导，其中 $ f_1(x) > 0 $, $ f_2(x) > 0 $, $ g(x) = e^{f_1(x)f_2(x)} $. 令 $ \varphi(x) = \frac{e^x g(x)}{f_1(x)f_2(x)} $，若 $ f_1'(0) = f_2'(0) = 0 $，则（ ）\\
    (A) $ \varphi'(0) = 0 $ \\
    (B) $ 0 < \varphi'(0) < 1 $ \\
    (C) $ \varphi'(0) > 1 $ \\
    (D) $ \varphi'(0) = 1 $

    \item 设椭圆 $ \frac{x^2}{a^2} + \frac{y^2}{b^2} = 1 $ 与抛物线 $ y = \frac{\sqrt{2}b}{a}(x - 1)^2 $ 相切于一点，记该切点的横坐标为 $ x_0 $，则（ ）\\
    (A) $ x_0 \in \left(0, \frac{1}{3}\right) $.\\
    (B) $ x_0 \in \left(\frac{1}{3}, \frac{2}{3}\right) $.\\
    (C) $ x_0 \in \left(\frac{2}{3}, 1\right) $.\\
    (D) 由已知条件不能确定 $ x_0 $ 的范围.

    \item 设函数 $ f(x) $ 有二阶连续导数，x=0 是 $ f(x) $ 的驻点，若在 x=0 附近， $ f(x) $ 满足 $ \sin x f''(x) + \tan x f'(x) = \sqrt[3]{\cos x} - 1 $，则（）\\
    (A) x = 0 是 $ f(x) $ 的极大值点. (B) x = 0 是 $ f(x) $ 的极小值点.\\
    (C) x = 0 不是 $ f(x) $ 的极值点. (D) 不能确定 x = 0 是否为 $ f(x) $ 的极值点.

    \item 已知函数 $ f(x) $ 在 x=0 处存在二阶导数 $ f''(0) $. 若 $ \lim_{x\to0}\frac{\sin x - x + xf(x)}{x^3} = 0 $，则下列选项中，正确的是（）\\
    (A) x=0 是 $ f(x) $ 的极小值点.\\
    (B) x=0 是 $ f(x) $ 的极大值点.\\
    (C) 点 $ (0, f(0)) $ 是 $ y = f(x) $ 的拐点.\\
    (D) x = 0 不是 $ f(x) $ 的驻点.

    \item 设 k 为任意非零常数，函数 $ f(x) $ 在 $ [a,b] $ 上满足 $ f'' - kf' - f = 0 $，且 $ f(a) = f(b) = 0 $，则下列说法中，正确的是（）\\
    (A) $ f(x) $ 在 $ [a,b] $ 上先单调减少，再单调增加. (B) $ f(x) $ 在 $ [a,b] $ 上先单调增加，再单调减少.\\
    (C) $ f(x) $ 在 $ [a,b] $ 上恒为零.\\
    (D) $ f(x) $ 在 $ (a,b) $ 内恰有一个零点.

    \item 设函数 $ f(x)=\int_{0}^{x}(t^{2}-4t+3)e^{t^{2}}dt,x\in[0,3] $，则下列说法中，正确的是（）\\
    (A) $ f(x) $ 为单调函数.\\
    (B) 4e - 9 为 $ f(x) $ 的一个上界.\\
    (C) $ f(x) $ 的最小值为 0 .\\
    (D) $ f(x) $ 不存在最大值.

    \item 若 $ f(x) $ 为区间 I 上的连续函数，且 $ f(x) $ 的值域包含于 $ I, x_1, x_2 $ 为 I 中任意两个不同的点，则下列命题中，正确的是（）\\
    (A) 若在区间 I 上， $ f(x)<0,f''(x)>0 $，则 $ f^{2}\left(\frac{x_{1}+x_{2}}{2}\right)<\frac{f^{2}\left(x_{1}\right)+f^{2}\left(x_{2}\right)}{2} $\\
    (B) 若在区间 I 上， $ f'(x)<0,f''(x)>0 $，则 $ f^{2}\left(\frac{x_{1}+x_{2}}{2}\right)<\frac{f^{2}\left(x_{1}\right)+f^{2}\left(x_{2}\right)}{2} $\\
    (C) 若在区间 I 上， $ f(x)>0,f''(x)>0 $ ，则 $ f\left(f\left(\frac{x_{1}+x_{2}}{2}\right)\right)<\frac{f\left(f\left(x_{1}\right)\right)+f\left(f\left(x_{2}\right)\right)}{2} $\\
    (D) 若在区间 I 上， $ f'(x)>0,f''(x)>0 $ ，则 $ f\left(f\left(\frac{x_{1}+x_{2}}{2}\right)\right)<\frac{f\left(f\left(x_{1}\right)\right)+f\left(f\left(x_{2}\right)\right)}{2} $

    \item 设函数 $ f(x) $ 在 $ [0,2] $ 上二阶可导，且 $ f(0) + f(2) = 0 $，则（）\\
    (A) 若 $ f''(x) > 0 $，则 $ \int_{0}^{2} f(x) \, dx > 0 $．\\
    (B) 若 $ f''(x) > 0 $，则 $ \int_{0}^{2} f(x) \, dx < 0 $\\
    (C) 若 $ f'(x) > 0 $，则 $ \int_{0}^{2} f(x) \, dx > 0 $。\\
    (D) 若 $ f'(x) > 0 $，则 $ \int_0^2 f(x) \, dx < 0 $.

    \item 设函数 $ f(x) $ 可导，且 $ f'(x) > 0, g(x) = \int_{0}^{x} f(t) \, dt $ 。若 $ g(1) = 1, g(3) = 7 $ ，则 $ g(2) $ 的值可能为（）\\
    (A) 2. \quad (B) 3. \quad (C) 4. \quad (D) 5.

    \item 设函数 $ f(x) $ 在 $ (0, +\infty) $ 上具有二阶导数， $ f''(x) < 0 $，记 $ u_n = f(n) $，则数列 $ \{u_n\} $ 发散是 $ u_1 > u_2 $ 的（）\\
    (A) 充分必要条件.\\
    (B) 充分不必要条件.\\
    (C) 必要不充分条件.\\
    (D) 既不充分也不必要条件.

    \item 设函数 $ f(x) $ 具有二阶导数， $ f(0) > 0 $, $ f'(0) < 0 $，当 $ x \geq 0 $ 时， $ f''(x) < 0 $，则下列说法中，正确的是（）\\
    (A) $ f(-f(0)f'(0)) > 0 $.\\
    (B) $ f\left(-\frac{f(0)}{f'(0)}\right) > 0 $.\\
    (C) $ f(x) $ 在 $ [0, +\infty) $ 上只有一个零点.\\
    (D) $ f(x) $ 在 $ [0, +\infty) $ 上没有零点.

    \item 设函数 $ f(x) $ 满足 $ f''(x) + \ln(x+1)f'(x) + \arctan x f(x) = \int_0^x f(t) \, dt $，则下列说法中，错误的是（ ）\\
    (A) 若 $ f'(0) \neq 0 $，则 x = 0 不是 $ f(x) $ 的极值点.\\
    (B) 若 $ f'(0) \neq 0 $，则点 $ (0, f(0)) $ 是曲线 $ y = f(x) $ 的拐点.\\
    (C) 若 x=0 是 $ f(x) $ 的极值点，则点 $ \left(0,f(0)\right) $ 是曲线 $ y=f(x) $ 的拐点.\\
    (D) 若 x=0 是 $ f(x) $ 的极值点，则当 $ x \to 0 $ 时， $ f'(x) $ 是 $ x^{2} $ 的高阶无穷小.

    \item 设函数 $ f(x) $ 在 $ [0, +\infty) $ 上二阶可导，且 $ f(0) = 0 $ 。下列命题中，正确命题的个数为（）\\
    ① 若在 $ (0,+\infty) $ 上恒有 $ f''(x)<0 $ ，则 $ \frac{f(x)}{x} $ 在 $ (0,+\infty) $ 上单调减少.\\
    ② 若 $ \frac{f(x)}{x} $ 在 $ (0,1] $ 上单调减少，则在 $ (0,1] $ 上恒有 $ f''(x) \leq 0 $ .\\
    ③ 若 $ \lim_{x\to+\infty}\frac{f(x)}{x} $ 存在，则 $ \lim_{x\to+\infty}f'(x) $ 存在.\\
    ④ 若在 $ (0,+\infty) $ 上恒有 $ f''(x)<0 $ ，且 $ \lim_{x\to+\infty}\frac{f(x)}{x} $ 存在，则 $ \lim_{x\to+\infty}f'(x) $ 存在.\\
    (A) 1. \quad (B) 2. \quad (C) 3. \quad (D) 4.

    \item 设函数 $ f(x) = \begin{cases} \frac{\sin x}{x}, & x \neq 0, \\ 1, & x = 0, \end{cases} $ 则下列区间中，是函数 $ g(x) = \int_0^x f(t) \, dt - \frac{2x}{\pi} $ 的单调增加区间为（ ）\\
    (A) $ [- \pi, 0] $ . (B) $ [0, \pi] $ . (C) $ \left[-\frac{\pi}{2}, \frac{\pi}{2}\right] $ . (D) $ [- \pi, \pi] $ .

    \item 设函数 $ f(x) $ 在 $ \left[0, \frac{\pi}{2}\right] $ 上可导，已知 $ f(0) = 0, f'(0) = 1 $，且 $ f(x) $ 满足 $ f'(x) \tan x \leq f(x) $，则下列关系式中，正确的是（）\\
    (A) $ f\left(\frac{\pi}{6}\right) \geq \frac{1}{2} $. \quad (B) $ f\left(\frac{\pi}{4}\right) \leq \frac{1}{\sqrt{2}} $. \quad (C) $ f\left(\frac{\pi}{4}\right) \leq \frac{\sqrt{2}}{2} f\left(\frac{\pi}{2}\right) $. \quad (D) $ f\left(\frac{\pi}{2}\right) \geq 2f\left(\frac{\pi}{6}\right) $.
\end{enumerate}

\subsubsection*{二、填空题}

\begin{enumerate}[start=21]
    \item 设函数 $ f(x)=\sqrt{\frac{1+x}{1-x}}\sqrt[4]{\frac{1+2x}{1-2x}} $ ，则 $ f'(0)= $ \underline{\hspace{2cm}}.

    \item 设 $ f(x) $ 为定义在 $ (-\infty,0) \cup (0,+\infty) $ 上的分段连续函数，且 $ \lim_{x\to0^-}f(x)=-1,\lim_{x\to0^+}f(x)=1 $，则 $ F(x)=\int_{0}^{x}(\sin x-\sin t)f(t)dt $ 在 x=0 处可导的最高阶数为 \underline{\hspace{2cm}}.

    \item 设函数 $ y = y(x) $ 由方程 $ \sin x - \int_{1}^{\ln y} e^{-t^{2}} \, dt = 0 $ 确定，则 $ y'(0) = $ \underline{\hspace{2cm}}.

    \item 设函数 $ y = y(x) $ 由方程 $ x + 1 = \int_{1}^{y+x} \sin^{2} t^{2} dt $ 确定，则 $ y''(-1) = $ \underline{\hspace{2cm}}.

    \item 设函数 $ y = y(x) $ 由方程组 $ \begin{cases} x = \arctan t, \\ \int_{0}^{2y} e^u \, du + \int_{t}^{0} \frac{\cos u}{1 + u^2} \, du = 0 \end{cases} $ 确定，则 $ y'(0) = $ \underline{\hspace{2cm}}.

    \item 已知函数 $ f(x) $ 在 $ (-\infty, +\infty) $ 上连续，且 $ f(x) = x^{2} + \int_{0}^{x}(x-t)f(t)dt $，则对任意正整数 $ n, f^{(2n)}(0) = $ \underline{\hspace{2cm}}.

    \item 设函数 $ f(x)=\left(x^{3}-1\right)^{n}\sin x $，则对任意正整数 $ n,f^{(n)}(1)= $ \underline{\hspace{2cm}}.

    \item 设函数 $ f(x) = \frac{\left[ax^{2} + (1 - 2a)x - 2\right]^{n}}{e^{x - 2}} $ ( $ a \neq 0 $)．若 $ f^{(n)}(2) = 2^{n}n $！，则 $ a = $ \underline{\hspace{2cm}}.

    \item 设函数 $ f(x) = x \ln x + \frac{e}{6} x^3 - \frac{e}{2} x^2 + \left( \frac{e}{2} - 1 \right) x - \frac{e}{6} $, $ g(x) = x (\ln x - 1) + e^x - ex $. 若曲线 $ y = f(x) $ 与 $ y = g(x) $ 相切于一点, 则它们的公共切线的方程为 \underline{\hspace{2cm}}.

    \item 已知质点以恒定速率沿曲线运动时，向心加速度大小为 $ \frac{v^{2}}{\rho} $，其中 v 为质点的运动速率， $ \rho $ 为曲率半径。设单位质点以恒定的速率 $ v_{0} $ 沿着一个形状满足 $ y = x^{2} $ 的光滑轨道移动，则其在点 $ (0,0) $ 处受到的向心力大小为 \underline{\hspace{2cm}}.

    \item 设 $ f(x) $ 为可微函数，曲线 $ y = f(x) $ 与曲线 $ y = \arctan x $ 在点 $ \left(1, \frac{\pi}{4}\right) $ 处有公共切线，则 $ \lim_{n \to \infty} n \left[ f\left( \frac{n}{n+1} \right) - \arctan \frac{n}{n-1} \right] = $ \underline{\hspace{2cm}}.

    \item 设函数 $ f(x) $ 具有二阶导数，且在点 $ (0,0) $ 处的曲率圆为 $ (x+1)^{2}+(y+2)^{2}=5 $，则 $ \lim_{n\to\infty}n^{2}\left[f\left(\frac{1}{n}\right)+f\left(-\frac{1}{n}\right)\right]= $ \underline{\hspace{2cm}}.

    \item 设函数 $ f(x) $ 具有二阶连续导数．若曲线 $ y = f(x) $ 过点 $ (0,0) $ 且与曲线 $ y = \ln x $ 在点 $ (e,1) $ 处相切，则 $ \int_{0}^{e} x f''(x) \, dx = $ \underline{\hspace{2cm}}.

    \item 曲线 $ x^{3}+y^{3}=y^{2} $ 的斜渐近线方程为 \underline{\hspace{2cm}}.

    \item 设函数 $ y = f(x) $ 在区间 $ [0, c) $ 上具有一阶连续导数，其值域为 $ [1, +\infty) $， $ f(0) = 1 $，并且 $ y' > y^2 $ 恒成立，则区间 $ [0, c) $ 内包含的最大整数是 \underline{\hspace{2cm}}.

    \item 设函数 $ f(x) $ 的定义域为 $ [0,c) $，其中 $ c>1 $，值域为 $ [1,+\infty) $， $ f(0)=1 $， $ f'(x)=\left|f(x)\right|^a $，则 $ a $ 的取值范围是 \underline{\hspace{2cm}}.
\end{enumerate}

\subsubsection*{三、解答题}

\begin{enumerate}[start=37]
    \item 设函数 $ f(x) $ 具有一阶连续导数，且 $ f(0)=1, f(1)=a $\\
    (I) 求使得 $ 1 + \frac{a}{\sqrt{2}} - \int_{0}^{1} \sqrt{1 + \left[ f'(x) \right]^{2}} \, dx $ 取得最大值的 $ f(x) $ 的表达式.\\
    (II) 将 $ 1 + \frac{a}{\sqrt{2}} - \int_{0}^{1} \sqrt{1 + \left[ f'(x) \right]^2} \, dx $ 取得的最大值记为 $ g(a) $，当 $ a $ 为何值时， $ g(a) $ 取得最大值？请求出该最大值。

    \item 设函数 $ f(x) $ 在 $ [0, +\infty) $ 上一阶可导，且当 x > 0 时， $ f'(x) \geq \frac{1}{x} $， $ f(0) < 0 $ 。证明： $ f(x) $ 在 $ (0, +\infty) $ 内有且仅有一个零点。

    \item 设函数 $ f(x) $ 在 $ (-\infty, +\infty) $ 上二阶可导， $ f''(x) > 0 $， $ x = c $ 是 $ f(x) $ 的最小值点， $ f(c) < 0 $。证明： $ f(x) $ 在 $ (-\infty, +\infty) $ 上恰有两个零点。

    \item 设函数 $ f(x) $ 在 $ [a,b] $ 上连续， $ g(x) $ 在 $ [a,b] $ 上有连续导数，且 $ g'(x) \neq 0 $ 。若 $ \int_{a}^{b} f(x) \, dx = 0, \int_{a}^{b} f(x) g(x) \, dx = 0 $ ，证明：至少存在两个不同的点 $ \xi_1, \xi_2 \in (a,b) $ ，使得 $ f(\xi_1) = f(\xi_2) = 0 $ 。

    \item 证明：当 x > 0 时， $ \left(x - e^{2}\right)\left(x \ln x - x - \frac{4}{3e}x \sqrt{x} + \frac{1}{3}e^{2}\right) \leq 0 $

    \item 证明: 当 x > 0 时, $ \frac{1}{2}x^{2}\ln^{2}x - \frac{3}{4}x^{2}\ln x + x\ln x - \frac{1}{4}\ln x \geq 0 $, 并指明等号成立的条件.

    \item 证明：对任意的 0 < a < 1, 0 < b < 1，均有 $ e^{\frac{-(a+b)^{2}}{8}} < e^{-\frac{a^{2}}{8}} + e^{-\frac{b^{2}}{8}} - 1 $。

    \item 设 $ f(x) $ 为 $ \left[-\frac{\pi}{2}, \frac{\pi}{2}\right] $ 上的连续函数，满足 $ f(-x) = f(x) $，且当 $ x \in \left(0, \frac{\pi}{2}\right] $ 时， $ f(x) = \frac{\sin x - x \cos x}{x(1 - \cos x)} $。证明： $ \frac{1}{2} \leq f(x) \leq \frac{2}{3} $。

    \item 已知函数 $ f(x) $ 在 $ \left[0, \frac{\pi}{2}\right] $ 上连续，在 $ \left(0, \frac{\pi}{2}\right) $ 内二阶可导。当 $ x \in \left(0, \frac{\pi}{2}\right) $ 时， $ f'(x) = \frac{x}{\sin x} + \frac{\sin x}{x} $，且 $ f(0) = 0 $。证明：对任意 $ a \in \left(0, \frac{\pi}{2}\right) $， $ af\left(\frac{a}{2}\right) < \int_{0}^{a} f(x) \, dx < \frac{a}{2}f(a) $。

    \item 证明： $ \frac{\sqrt{2}\pi^{2}}{8}\leq\int_{0}^{+\infty}\frac{\arctan x}{x\sqrt{1+x^{2}}}\mathrm{d}x\leq\frac{\pi(\pi+2)}{8} $

    \item 设函数 $ f(x) $ 在 $ [0,3] $ 上连续，在 $ (0,3) $ 内可导，且 $ f(0) + 2f(1) = 2f(2) + f(3) = 1 $.\\
    证明：存在 $ \xi \in (0,3) $，使得 $ f'(\xi) = 0 $。

    \item 设函数 $ f(x) $ 在 $ [1,e^2] $ 上连续，在 $ (1,e^2) $ 内具有二阶导数，且 $ f(1)=0 $， $ f(e)=1 $， $ f(e^2) $ =2. 证明：存在 $ \xi \in (1,e^2) $，使得 $ f''(\xi) \xi^2 = -1 $.

    \item 设函数 $ f(x) $ 在 $ \left[0, \frac{\pi}{2}\right] $ 上二阶可导，且 $ f(0) = 0 $， $ f\left(\frac{\pi}{2}\right) = f'(0) = 1 $ 。证明：存在 $ \xi \in \left(0, \frac{\pi}{2}\right) $，使得 $ f''(\xi) = -\sin\xi $ 。

    \item 设 $ [0,1] $ 上的连续函数 $ f(x) $ 满足 $ \lim_{x\to0^+}\frac{f(x)}{x}=\lim_{x\to1^-}\frac{f(x)-1}{x-1}=-1,f\left(\frac{1}{2}\right)=\frac{1}{4} $．证明：存在 $ \xi\in(0,1) $，使得 $ f''(\xi)=2 $．

    \item 设函数 $ f(x) $ 在 $ x_0 $ 的一个去心邻域 $ U(x_0) $ 内二阶可导，证明：对 $ \hat{U}(x_0) $ 内的任意一点 $ x $，都存在一点 $ \xi_x $，使得 $ f(x) - f(x_0) = f'(x_0)(x - x_0) + f''(\xi_x)(x - \xi_x)(x - x_0) $，其中 $ \xi $ 介于 $ x $ 与 $ x_0 $ 之间。

    \item 设函数 $ f(x) $ 在 $ [0,1] $ 上二阶可导，且 $ f(0) = f'(0) = 1 $ 。证明：对任意 $ x \in (0,1) $，均存在 $ \xi_x \in (0,1) $，使得 $ f(x) - x - 1 = \frac{x^2}{2} f''(\xi_x) $ 。

    \item 设 $a>0$，函数 $f(x)$ 在 $[a,b]$ 上连续，在 $(a,b)$ 内可导。证明：若对任意 $x\in(a,b)$，都有 $f'(x)\neq0$，则存在 $\xi_1,\xi_2\in(a,b)$，使得 $\frac{f'(\xi_1)}{f'(\xi_2)}=\frac{a^2+ab+b^2}{a+b}\cdot\frac{2\xi_1}{3\xi_2}^2$。

    \item 设函数 $ f(x) $ 在 $ [0,e] $ 上连续，在 $ (0,e) $ 内可导，且 $ f(0)=0 $ 。证明：存在 $ \xi, \eta \in (0,e) $，使得 $ \xi f(\xi) + (\xi^2 - \xi) f'(\xi) = (e^2 - e) f'(\eta) $。

    \item 证明: 在区间 $ (0,2) $ 内存在三个不同的点 $ x_{1}, x_{2}, x_{3} $，使得 $ \frac{1-\ln(1+x_{1})}{(1+x_{1})^{2}}x_{3}=\frac{1-\ln(1+x_{2})}{(1+x_{2})^{2}}(2-x_{3}) $.

    \item 设函数 $ f(x) $ 在 $ [0,1] $ 上二阶可导， $ f(0)=0 $， $ f(1)=1 $，且 $ f''(x)<0 $ 。证明：存在唯一的 $ \xi \in (0,1) $，使得 $ f(\xi)=1-\xi $，且 $ \xi \leq \frac{1}{2} $ 。

    \item 设函数 $ f(x) $ 在 $ [0,1] $ 上三阶可导，且 $ f(1) = f'(1) = f''(1) = 2f(0) $ 。证明：存在 $ \xi \in (0,1) $，使得 $ f'''(\xi) = 0 $ 。

    \item 设函数 $ f(x) $ 为 $ [0,\pi] $ 上的连续正值函数， $ F(x) $ 为 $ [0,\pi] $ 上的连续函数，且对 $ x \in (0,\pi) $， $ F(x) = \frac{\int_{0}^{x}(1 - \cos t)f(t)dt}{\int_{0}^{x}t^{2}f(t)dt} $。求 $ F(x) $ 在 $ [0,\pi] $ 上的最大值。

    \item 设函数 $ f(x) $ 为 $ [1,+\infty) $ 上的正值函数，且 $ f(x) \leq e^{\sqrt{x}} $ 。证明：存在 $ \xi \in (1,+\infty) $，使得 $ f'(\xi) < f(\xi) $。

    \item 设函数 $ f(x) $ 为 $ [0, +\infty) $ 上恒正且具有二阶连续导数的凹函数，且 $ f'(1) > 0 $ 。证明：存在唯一的 $ x_0 \in (0, 1) $，使得曲线 $ y = f(x) $ 在点 $ (x_0, f(x_0)) $ 处的切线在 y 轴上的截距等于 $ [0, x_0] $ 上以 $ y = f(x) $ 为曲边的曲边梯形的面积。

    \item 设函数 $ f(x) $, $ g(x) $ 在 $ [a,b] $ 上连续，在 $ (a,b) $ 内具有二阶导数，且 $ f(a) = g(b) $, $ f(b) = g(a) $, $ \int_{a}^{b} f(x) \, dx = \int_{a}^{b} g(x) \, dx $. 证明: 存在 $ \xi, \eta \in (a,b) $，使得 $ f''(\xi) = g''(\eta) $.

    \item 设函数 $ f(x) $ 满足 $ f'(x) = \frac{1}{x \ln^{2} x + f^{2}(x)} (x \geq e) $，且 $ f(e) = 0 $，证明：\\
    (I) 当 x > e 时， $ f'(x) < \frac{1}{x \ln^{2} x} $;\\
    (II) $ \lim_{x\to+\infty}f(x) $ 存在, 且其值严格小于 1 .
\end{enumerate}

\subsection{c 类}

\subsubsection*{一、选择题}

\begin{enumerate}
    \item 已知函数 $ f(x)=\begin{cases}\sin x, & x\in Q\cap(-1,1), \\ \tan x, & x\in(-1,1)\setminus Q,\end{cases} $ 其中 Q 为有理数，则（）\\
    (A) $ f(x) $ 在 x=0 处可导.\\
    (B) $ f(x) $ 在 x=0 处连续但不可导.\\
    (C) x=0 是 $ f(x) $ 的第一类间断点.\\
    (D) x=0 是 $ f(x) $ 的第二类间断点.
\end{enumerate}

\subsubsection*{二、填空题}

\begin{enumerate}[start=2]
    \item 设函数 $ y = y(x) $ 由参数方程 $ \begin{cases} x = \int_{0}^{u} e^{t^{2}} \, dt, \\ y = \int_{0}^{1} |e^{u} - e^{t}| \, dt \end{cases} $ 确定，若 $ \int_{0}^{\frac{1}{2}} e^{x^{2} - \frac{1}{4}} \, dx = a $，则 $ y_{+}^{'}(0) = $ \underline{\hspace{2cm}}.
\end{enumerate}

\subsubsection*{三、解答题}

\begin{enumerate}[start=3]
    \item 设函数 $ f(x) $ 在 $ [0,2] $ 上具有三阶连续导数，且 $ f(0)=1, f(1)=2, f(2)=4 $ 。证明：存在一个二次多项式 $ P(x)=ax^{2}+bx+c $ ，使得对任意 $ x \in [0,2] $ ，存在 $ \eta_{x} \in (0,2) $ ，满足
    $$ f(x)=P(x)+\frac{f''(\eta_{x})}{6}x(x-1)(x-2). $$

    \item 设函数 $ f(x) $ 二阶可导，且 $ f''(0) \neq 0 $， $ a, b $ 为两个不相等的实数。若存在函数 $ \theta(x) $，使得 $ f(ax) - f(bx) = f'(\theta(x)x)(a - b)x $， $ \theta(x) $ 介于 $ a, b $ 之间，求 $ \lim_{x \to 0} \theta(x) $。
\end{enumerate}

\end{document}